\documentclass[a4paper,11pt,reqno]{amsart}

\usepackage[textwidth=31pc,textheight=48pc,centering]{geometry}
\usepackage[T1]{fontenc}
\usepackage{lmodern}
\usepackage{microtype}
\usepackage{etoolbox}
\usepackage{amsmath,amssymb,mathtools}
\usepackage{mathrsfs}
\usepackage[numbers,sort&compress]{natbib}
\usepackage{enumitem}
\usepackage{xurl}
\usepackage[hidelinks,pdfencoding=auto,psdextra]{hyperref}
\usepackage{bookmark}

% Annals-like text dimensions; standard AMS author-review manuscript.
% No journal banner, publication DOI, acceptance date or publisher copyright.
\patchcmd{\abstract}{\leftmargin3pc}{\leftmargin1.5pc}{}%
  {\ClassError{amsart}{Abstract margin patch failed}{Restore the standard abstract environment.}}

% Author-review layout: metadata belongs with the abstract, not in footnotes
% below the first theorem. The ordinary AMS metadata commands are retained.
\makeatletter
\def\@adminfootnotes{%
  \let\@makefnmark\relax \let\@thefnmark\relax
  \ifx\@empty\thankses\else \@footnotetext{%
    \def\par{\let\par\@par}\@setthanks}%
  \fi}
\newcommand{\paper@frontmatter}{%
  \par\smallskip\begingroup\small
  \setlength{\parindent}{0pt}%
  \setlength{\parskip}{0.25\baselineskip}%
  \textit{2020 Mathematics Subject Classification.}\enspace\@subjclass.\par
  \textit{Keywords.}\enspace\@keywords.\par
  \endgroup}
\patchcmd{\@maketitle}{\@setabstract}%
  {\begin{center}\small\@date\end{center}\@setabstract\paper@frontmatter}{}%
  {\ClassError{amsart}{Front-matter placement patch failed}{Restore the standard AMS title block.}}
\makeatother

\setcounter{secnumdepth}{1}
\setcounter{tocdepth}{1}
\setlength{\parindent}{1.5em}
\setlength{\parskip}{0pt}
\setlength{\emergencystretch}{3em}
\setlist[itemize]{leftmargin=2em,itemsep=0.12em,topsep=0.35em}
\setlist[enumerate]{leftmargin=2.2em,itemsep=0.12em,topsep=0.35em}
\allowdisplaybreaks[2]
\raggedbottom
\Urlmuskip=0mu plus 1mu
\frenchspacing
\tolerance=1800
\pretolerance=500
\ifdefined\pdfgentounicode
  \pdfgentounicode=1
\fi

\providecommand{\tightlist}{\setlength{\itemsep}{0pt}\setlength{\parskip}{0pt}}
\providecommand{\passthrough}[1]{#1}
\newcommand{\proofpart}[1]{%
  \par\smallskip\noindent\textit{#1.}\enspace\ignorespaces}

\theoremstyle{plain}
\newtheorem*{theoremA}{Theorem A}
\newtheorem{theorem}{Theorem}[section]
\newtheorem{lemma}[theorem]{Lemma}
\newtheorem{proposition}[theorem]{Proposition}
\newtheorem{corollary}[theorem]{Corollary}
\newtheorem*{claim}{Claim}
\theoremstyle{definition}
\newtheorem{definition}[theorem]{Definition}
\newtheorem{problem}[theorem]{Problem}
\theoremstyle{remark}
\newtheorem{remark}[theorem]{Remark}

\hypersetup{
  hidelinks,pdfencoding=auto,psdextra,
  pdftitle={Obligatory Triple Systems: Canonical Atoms and Exact Finite Spectra},
  pdfauthor={Samuil Petkov},
  pdfsubject={A canonical atom decomposition, exact finite spectra, and Lean verification for obligatory triple systems},
  pdfkeywords={Erdos Problem 593, obligatory triple systems, uncountable chromatic number, canonical decomposition, partition lattices, Lean 4},
  pdfcreator={LaTeX with amsart; author-review manuscript},
  pdfproducer={pdfTeX},
  pdfdisplaydoctitle=true
}

\title[Obligatory triple systems]{Obligatory triple systems: canonical atoms and exact finite spectra}
\author{Samuil Petkov}
\address{École normale supérieure--PSL, 45 rue d'Ulm, 75005 Paris, France}
\email{samuil.petkov@ens.psl.eu}
\subjclass[2020]{Primary 05C65; Secondary 05C15, 05C35, 05C40, 05C63, 06A07, 03E05}
\keywords{Erdős Problem 593, obligatory triple systems, uncountable chromatic number, canonical decomposition, partition lattices, entropy, homomorphism densities, Lean~4}
\date{9 October 2026}

\begin{document}

\begin{abstract}
Erdős Problem~593 asks which finite triple systems occur in every triple
system of uncountable chromatic number.  We give an independent proof of the
classification and develop its finite structural consequences.  After
isolated vertices are removed, a system is obligatory if and only if it is
linear, every hyperedge-node of its Levi graph is incident with a bridge,
and every Berge cycle has even length.  Its canonical atoms are single
triples and private-vertex expansions of finite $2$-connected bipartite
graphs.  For reduced obligatory systems with at least one hyperedge, we
determine the exact parameter and atom-count spectra, including the parity
obstruction at cycle rank one.  The supported one-point factorizations form
a product of partition lattices: every attachment excess profile can be
realized with the atoms and component count fixed, and the lattice type
recovers that profile.  For connected systems of positive cycle rank, an
exact budget relates cyclic-rank splitting and excess core order to the
available factorization rank.  At every prescribed original-edge piece
count we obtain sharp binomial and Stirling-number bounds, with binary
joints and a common joint as the interior equality cases.  The two
extremal assemblies can have the same atoms and global parameters.
These results separate the constraints imposed by cyclic content from
the freedom of assembly.  For the associated graph cores, we also prove
exact edge-marginal entropy bounds for uniform Potts laws and nearby general
pair laws, and a common-power density bound for regular hosts with two
centered spectral levels.  These quantitative comparisons are separate
from the obligatory classification; their unrestricted versions remain open.
The classification and specified structural
extensions are verified in Lean~4; profile realization, full lattice
spectra, the new counting endpoints and the entropy--density arguments
are not yet fully formalized.
All arguments are in ZFC.
\end{abstract}

\maketitle

\section*{Introduction}\label{introduction}

Although the host systems are infinite, the classification is determined by
finite structure.  A finite triple system is \emph{obligatory} if it occurs
in every triple system of uncountable chromatic number.  Erdős asked for a
characterisation of these finite systems \citep{erdos1995}; the question is
catalogued as Erdős Problem~593 \citep{bloom593}.  We give both constructive
and intrinsic descriptions of the answer.

For a finite graph $J$, write $J^+$ for its private-vertex expansion.  Let
$\mathcal B$ be the smallest class containing $J^+$ for every finite
bipartite graph $J$ and every finite edgeless triple system, and closed under
finite disjoint unions and one-point amalgamations.  Membership in
$\mathcal B$ is understood up to triple-system isomorphism.  If $F$ is a
triple system, let $F^\circ$ denote the system obtained by deleting its
isolated vertices.

\begin{theoremA}\label{theorem-a-exact-characterisation}
For every finite triple system $F$, the following are equivalent:
\begin{enumerate}
\item $F$ is obligatory;
\item $F\in\mathcal B$;
\item $F^\circ$ is linear, every hyperedge-node of $I(F^\circ)$ is incident
with a bridge, and every Berge cycle of $F^\circ$ has even length.
\end{enumerate}
\end{theoremA}

The classification first appeared in Li's preprint, posted on 23 June 2026
\citep[Theorem~1.1]{li2026}.  The present proof and Lean implementation were
developed independently.  Both proofs use one-apex sequence lifts and Levi
bridges, but their finite trace arguments differ: Li uses a bridge-trace
formulation and graph derivatives \citep[Sections~3--4 and Theorem~4.6]{li2026},
whereas we decompose traces into expansion fibres and a support-incidence
forest.  We cite Li's work for chronology and comparison; none of its results
is used below as a black box.
\subsection*{Finite structure beyond the classification}
Theorem~A answers an embedding question about infinite hosts.  The further
results concern a different question: how much finite structure is fixed by
cyclic content, and how much remains free under one-point assembly?  The finite
systems considered in this overview are reduced (they have no isolated
vertices), obligatory, and have at least one hyperedge.  Write \(n=|V(F)|\),
\(m=|E(F)|\), and \(c\) for their number of components; put \(s=n-m\) and
\(\beta=2m-n+c\), the Levi cycle rank.  The
canonical normal form identifies the atoms as single triples and expansions
of finite \(2\)-connected bipartite graphs
(Theorem~\ref{theorem-canonical-atom-normal-form}).  Block decomposition itself
is classical \citep{whitney1932,diestel2017}; the point here is the exact list
of allowable block types and what it implies about the systems assembled
from them.

The finite results give two complementary classifications.  The numerical
spectra determine which global parameters and canonical atom counts are
possible.  The local partition product determines which decompositions
remain available once the atoms have been fixed.  Neither classification
alone captures the other: even assemblies with identical atoms, order,
size and cycle rank can have different factorization lattices and different
numbers of decompositions.

The numerical consequences are exact existence statements, not only bounds.
Theorem~\ref{theorem-order-size-component-spectrum} determines every feasible
order, size and component count.
Theorem~\ref{theorem-exact-canonical-atom-count-spectrum} and
Proposition~\ref{proposition-componentwise-atom-count-spectrum} then determine
all possible canonical atom counts \(k\) at fixed \(s,\beta,c\), including
the parity obstruction when \(\beta=1\).  For connected systems, the
indecomposability phase diagram
separates parameter ranges where every system is decomposable, where both
behaviours occur, and where every system is indecomposable
(Theorem~\ref{theorem-indecomposability-phase-diagram}).  Its boundary cases
identify cycle, even-theta and complete-bipartite expansions, and the
extremal decomposable assembly
(Corollary~\ref{corollary-structural-boundary-rigidity}).  These conclusions
do not say that numerical parameters determine every core up to isomorphism.

The remaining freedom is described by a lattice, not another numerical
invariant.  At each point \(p\) shared by \(\mu(p)\) canonical atoms, one may
choose a partition of those atoms independently.  The supported one-point
forest decompositions of the original hyperedge set, ordered by refinement,
therefore form \(\mathcal D(F)\cong\prod_{p\in S(F)}\Pi_{\mu(p)}\), where
\(S(F)\) is the set of shared points and \(\Pi_r\) is the partition lattice
on \(r\) elements (Lemma~\ref{lemma-local-decomposition-lattice-product}).
Its rank is \(N=k-c\).  Every integer partition \(\lambda\vdash N\) is
realized as an attachment excess profile by reassembling the same atoms
with the same component count, order, size and cycle rank
(Lemma~\ref{lemma-capacity-safe-profile-realization}).  This gives the exact
factorization-lattice spectrum: its types are
\(\prod_i\Pi_{\lambda_i+1}\), and the classical characteristic polynomials
of partition lattices \citep[Chapter~3]{stanley2012} recover \(\lambda\)
(Corollary~\ref{corollary-factorization-lattice-spectrum}).  The classification
is of lattice types,
not of all triple systems with the same parameters.

Cyclic content and assembly freedom are linked by an exact budget.  For a
connected system of positive cycle rank, the lattice rank \(N=k-1\), the
cost of distributing positive rank among several atoms, and the total
core-order excess above the fixed-rank minima sum to
\(s-2-\lceil2\sqrt{\beta}\rceil\)
(Proposition~\ref{proposition-cyclic-cost-budget}).  Thus rank splitting
and excess core order reduce the available decomposition rank.  Zero cost
means that one minimum-order cyclic atom carries all positive rank and
the remaining atoms are single triples; it does not imply a unique
isomorphism type for the cyclic core.

Assembly geometry also controls the number of decompositions at each
prescribed piece count.  If \(d_j(F)\) counts decompositions into exactly
\(k-j\) nonempty original-edge pieces, then
\(\binom Nj\le d_j(F)\le S(N+1,N+1-j)\) for \(0\le j\le N\), where
\(S(a,b)\) is a second-kind Stirling number \citep{comtet1974}.
For \(N\ge2\), equality at any interior \(j\) characterizes binary joints
at the lower bound and a single common joint at the upper bound
(Corollary~\ref{corollary-sharp-piece-counts}).  Both extremal profiles can
be attained with the atom list and component count fixed.  The example
following that corollary has identical atoms and global parameters but
four or five decomposition choices.  These counting consequences quantify
the freedom left by the canonical structure; they are not assertions of
new general Stirling-number identities.

\subsection*{Quantitative questions suggested by the atoms}
The description of the cyclic atoms also identifies graph cores for finite
quantitative questions.  Every simple bipartite subdivision of $K_4$ is
$2$-connected and has cycle rank three, so its private-vertex expansion is
an allowable atom.  This observation motivates, but does not prove, a density
comparison for the original graph.  Section~\ref{sec:quantitative-cores}
studies exact pair-marginal entropy witnesses and regular-host homomorphism
densities for these cores.  The first results cover uniform $q$-state Potts
tables, arbitrary tables in specified open families, and hosts with two
centered spectral levels.  They are not a universal Sidorenko theorem.

The entropy formulation has established predecessors
\citep{csikvarilin2017,szegedy2015}; local density comparisons near independence
are also classical \citep{lovasz2010}.  The positive-spectrum comparison has
its own prior literature \citep{sidorenko2021}.  We give the finite arguments
needed for our restricted statements and separate them from the two
unrestricted questions.  Neither the graph comparisons nor their novelty
is inferred from the classification of obligatory triple systems.
Some of these core graphs already belong to known scalar-density classes:
graphs with a vertex complete to the opposite part are covered in
\citet{conlonfoxsudakov2012}.  Such coverage is not a new graph family here.
More recently, \citet{imliliu2026} proved subdivision and even-theta
substitution results.  Their uniform substitution theorem and their
nonuniform theorem with path-count divisibility hypotheses do not settle
arbitrary choices of the six subdivision lengths considered below.

The classification and substantial structural extensions have accepted
Lean implementations.  The distinction is between the mathematical results
proved in the manuscript and the exact interfaces already exported by Lean,
not between a numerical test and a proof.  Full profile realization, lattice
spectra, maximal-chain and height endpoints, and the new counting statements
still have unfinished formal interfaces.  The finite results above have
manuscript proofs; they are not all claimed as accepted Lean exports.
The entropy and density arguments have no accepted Lean implementation.
Appendix~\ref{appendix-formalization} records the implementation routes,
validation records and remaining obligations without interrupting the
mathematical arguments.

\subsection*{Context and proof strategy}

The private-vertex expansion $J^+$ also appears in finite extremal
hypergraph theory.  Kostochka, Mubayi and Verstra{\"e}te study its
Tur{\'a}n extremal function \citep{kostochkamubayiverstraete2015}.
Here the same operation supplies the canonical cyclic pieces, but
the host condition is uncountable chromatic number rather than edge density.

The graph analogue was proved by Erdős and Hajnal: the obligatory finite
graphs are exactly the bipartite graphs
\citep[Corollary~5.6 and Theorem~7.4]{erdoshajnal1966}.  The classical
nonlinearity obstruction for uniform hypergraphs is due to Erdős, Hajnal,
and Rothschild \citep[Theorem~2, p.~532]{ehr1973}.  Obligatory triple
systems were studied further by \citet{egh1975}, \citet{komjath2001},
\citet{hajnalkomjath2008}, and \citet{komjath2008}; the expansion theorem of
\citet[Theorem~1.2]{reiher2024} supplies the positive atoms used below.

In the finite setting, arbitrarily large chromatic number is compatible with
arbitrarily large girth: Erd\H{o}s proved this for graphs
\citep[pp.~35--37]{erdos1959}, and Erd\H{o}s and Lov\'asz established a
uniform-hypergraph counterpart \citep[Theorem~1$'$]{erdoslovasz1975}.
Later constructions control sparsity quantitatively: Kostochka and R{\"o}dl
obtain high-girth uniform hypergraphs with controlled maximum degree
\citep[Theorem~3]{kostochkarodl2010}, while Alon, Kostochka, Reiniger, West
and Zhu give an explicit construction \citep[Theorem~3.2]{alonetal2016}.
Uncountable chromatic number is essential here.  For finite hosts,
\citet[Proposition~1.3(ii)]{wang2026} record that the weak chromatic numbers
of $F$-free uniform hypergraphs are bounded exactly when $F$ contains no Berge
cycle.  This does not classify obligatory systems: the expansion $C_4^+$
contains a Berge cycle and can be avoided by finite triple systems of
arbitrarily large chromatic number, yet it occurs in every uncountably
chromatic triple system by Theorem~A.  Finite high-chromatic avoidance and
uncountable obligatoriness are therefore different questions.

The proof separates into positive and negative directions.  For the positive
direction, a probabilistic rainbow lemma provides the local injectivity
needed to force complete bipartite expansions.  A rooted-abundance lemma then
establishes closure under one-point amalgamation, including at singular
uncountable cardinals.  After the bridges of the Levi graph are deleted, each
remaining active component is the expansion of a finite bipartite graph.  The
quotient is a forest, whose running-intersection property determines a valid
amalgamation order.

For the negative direction, we use an independently developed one-apex
sequence lift; \citet[Sections~3--4]{li2026} gives a parallel formulation.
The lift transfers uncountable chromatic
number from a graph to a triple system while controlling every finite linear
trace.  The trace theorem is proved by decomposing the trace into expansion
fibres whose support-incidence graph is a forest, so distinct fibres can meet
only at cut points.  Consequently, neither a bridge-free hyperedge-node nor a
Berge cycle can be assembled across fibres.  Each of the three failures of
the intrinsic criterion has a corresponding host: an Erdős--Rado linear host
excludes nonlinearity, the lift of $K_{\omega_1}$ excludes a missing bridge,
and a classical Erdős--Hajnal graph of uncountable chromatic number whose odd
girth exceeds a prescribed finite bound excludes an odd Berge cycle.

All embeddings are injective and non-induced, and the argument is carried out
in ZFC.  For the classification proof, the imported black-box inputs are the de Bruijn--Erdős
compactness theorem \citep[Theorem~1, pp.~371--373]{debruijn1951}, the
Erdős--Hajnal high-odd-girth theorem
\citep[Theorem~7.4, p.~76]{erdoshajnal1966}, and the pair case of the
Erdős--Rado partition theorem
\citep[Theorem~4(i), formula~(95), p.~471]{erdosrado1956}.
The later block-formulation corollary also uses Komjáth's block reduction
\citep{komjath2001}.
For general background on cardinal partition relations, see
\citet{ehmr1984}; the specific imported statements are those listed above.

\section{Preliminaries}\label{conventions-and-definitions}

A \emph{triple system} is a simple \(3\)-uniform hypergraph \(H=(V(H),E(H))\), so \(E(H)\subseteq[V(H)]^3\). A colouring \(c:V(H)\to\lambda\) is \emph{proper} when no hyperedge is monochromatic, and \(\chi(H)\) is the least cardinal \(\lambda\) admitting such a colouring. An embedding of a triple system \(F\) in \(H\) is an injective map \(f:V(F)\to V(H)\) such that \(f[e]\in E(H)\) for every \(e\in E(F)\); embeddings are not required to be induced. A finite triple system \(F\) is \emph{obligatory} if every triple system \(H\) with \(\chi(H)>\aleph_0\) contains an embedding of \(F\). A point is \emph{isolated} if it belongs to no hyperedge, and \(F^\circ\) denotes the result of deleting all isolated points.

All ordinary graphs in this paper are simple. A triple system is \emph{linear} when any two distinct hyperedges meet in at most one point. Its \emph{Levi graph} \(I(F)\) is the bipartite graph with classes \(V(F)\) and \(E(F)\), where a point-node \(p\) is adjacent to a hyperedge-node \(e\) exactly when \(p\in e\). A \emph{bridge} is an edge of a graph whose deletion increases the number of connected components. A \emph{Berge cycle of length} \(\ell\ge2\) consists of distinct points \(p_0,\ldots,p_{\ell-1}\) and distinct hyperedges \(e_0,\ldots,e_{\ell-1}\) such that \(p_i\in e_i\cap e_{i+1}\), with indices modulo \(\ell\). Equivalently, it is a simple cycle of length \(2\ell\) in \(I(F)\). For graph terminology and the incidence-graph interpretation of hypergraph connectivity, see \citet[Chapters~1 and~3]{diestel2017} and \citet{bahmanian2015}.

For a finite simple graph \(J\), its \emph{private-vertex expansion} \(J^+\) has vertex set
\[
V(J)\mathbin{\dot\cup}\{p_a:a\in E(J)\}
\]
and hyperedges \(\{u,v,p_{\{u,v\}}\}\) for \(\{u,v\}\in E(J)\), where the points \(p_a\) are new and pairwise distinct. A \emph{one-point amalgamation} of vertex-disjoint triple systems \(F_0,F_1\) is obtained by choosing \(x_i\in V(F_i)\), identifying \(x_0\) with \(x_1\), and making no other identifications or new hyperedges. Finite disjoint unions and finite edgeless systems include the empty cases.

\paragraph{Finite reductions.}\label{finite-reductions}

\begin{lemma}[Isolated-vertex reduction]\label{lemma-1.1-isolated-vertex-reduction}

Let \(F\) be a finite triple system and let \(F^\circ\) be obtained by deleting all isolated vertices. Then

\[
F\text{ is obligatory}\iff F^\circ\text{ is obligatory},
\]

and

\[
F\in\mathcal B\iff F^\circ\in\mathcal B.
\]

\end{lemma}
\begin{proof}
If \(F^\circ\) is obligatory and \(H\) has uncountable chromatic number, then \(H\) has infinitely many vertices. After embedding the finite system \(F^\circ\), map the isolated vertices of \(F\) injectively to unused host vertices. Containment is non-induced, so additional host edges are irrelevant. The converse follows because \(F^\circ\) is a subhypergraph of \(F\).

For membership in \(\mathcal B\), the implication \(F^\circ\in\mathcal B\Rightarrow F\in\mathcal B\) follows by taking the disjoint union with a finite edgeless system.

For the converse, use structural induction on a construction of \(F\) in \(\mathcal B\). If \(F=J^+\), delete the isolated graph vertices of \(J\) to obtain \(J_0\); then \((J^+)^\circ=J_0^+\in\mathcal B\). If \(F\) is edgeless, its reduction is the empty edgeless system. Reduction commutes with disjoint union.

Suppose finally that \(F\) is a one-point amalgamation of \(F_0\) and
\(F_1\), identifying \(x_i\in V(F_i)\). By induction,
\(F_0^\circ,F_1^\circ\in\mathcal B\). If both selected vertices are
nonisolated in their factors, then \(F^\circ\) is the one-point amalgamation
of the two reductions at those vertices. If exactly one selected vertex is
nonisolated, the isolated side disappears at the identified point and
\(F^\circ\) is the disjoint union of the two reductions. The same description
applies when both selected vertices are isolated, because the identified
point is then deleted. Hence \(F^\circ\in\mathcal B\) in every case.
\end{proof}

We may therefore suppress isolated vertices in the substantive argument and
restore them at the end.

\begin{lemma}[Finite deletion]\label{lemma-1.2-finite-deletion}

If \(\chi(H)>\aleph_0\) and \(S\subseteq V(H)\) is finite, then

\[
\chi(H-S)>\aleph_0.
\]

\end{lemma}
\begin{proof}
Otherwise a countable colouring of \(H-S\), together with finitely many fresh colours for \(S\), would colour \(H\) countably.
\end{proof}

\paragraph{Colouring and closure tools.}\label{colouring-and-closure-tools}

\begin{lemma}[Two elementary colouring facts]\label{lemma-1.3-two-elementary-colouring-facts}

We use the following two facts repeatedly.

\begin{enumerate}
\def\labelenumi{\arabic{enumi}.}
\tightlist
\item
  If \(E(H)=E_1\cup\cdots\cup E_r\) and every \((V(H),E_i)\) is countably chromatic, then \(H\) is countably chromatic.
\item
  If a graph \(G\) admits an orientation in which every vertex has outdegree at most \(d<\omega\), then \(\chi(G)\le 2d+1\).
\end{enumerate}

\end{lemma}
\begin{proof}
For the first assertion, take the product of the \(r\) countable colourings.

For the second, every finite subgraph has at most \(d|V|\) edges and therefore a vertex of degree at most \(2d\). It is \((2d+1)\)-colourable by greedy deletion. The de Bruijn--Erdős compactness theorem \citep[Theorem~1]{debruijn1951} then gives a \((2d+1)\)-colouring of all of \(G\).
\end{proof}

\begin{lemma}[Closure-chain lemma]\label{lemma-1.4-closure-chain-lemma}

Let \(\kappa\) be an uncountable cardinal, let \(1\le r<\omega\), and let

\[
\Phi:[\kappa]^r\longrightarrow[\kappa]^{<\omega}
\]

have uniformly finite values. There is an increasing continuous chain

\[
\varnothing=M_0\subseteq M_1\subseteq\cdots
\subseteq M_i\subseteq\cdots\qquad(i<\operatorname{cf}\kappa)
\]

such that

\[
\bigcup_{i<\operatorname{cf}\kappa}M_i=\kappa,
\qquad |M_i|<\kappa,
\]

and every \(M_i\) is \(\Phi\)-closed:

\[
a\in[M_i]^r\quad\Longrightarrow\quad\Phi(a)\subseteq M_i.
\]

\end{lemma}
\begin{proof}
Choose sets \(A_i\subseteq\kappa\), each of cardinality below \(\kappa\), whose union is \(\kappa\). At a successor stage close \(M_i\cup A_i\) under \(\Phi\), iterating the operation \(\omega\) times. Since \(\Phi\) is finitary and finite-valued, this does not increase an infinite cardinal below \(\kappa\). At a limit stage take the union. A union of an increasing chain of sets closed under a finitary operation remains closed; and because the stage is below \(\operatorname{cf}\kappa\), its cardinality remains below \(\kappa\).
\end{proof}

Write

\[
I_i=M_{i+1}\setminus M_i.
\]

Then \((I_i)_{i<\operatorname{cf}\kappa}\) partitions \(\kappa\), and \(M_i=\bigcup_{j<i}I_j\).

\section{A bipartite subgraph lemma}\label{the-graph-lemma-needed-for-the-expansion-atom}

The following finite complete-bipartite lemma is a consequence of
\citet[Corollary~5.6, p.~72]{erdoshajnal1966}.  We give a short
closure-chain proof because the construction will also be used later.

\begin{lemma}[Uncountable chromatic number forces complete bipartite graphs]\label{lemma-2.1-uncountable-chromatic-number-forces-complete-bipartite-graphs}

For every positive integer \(n\), every graph of uncountable chromatic number contains \(K_{n,n}\).

\end{lemma}
\begin{proof}
Suppose otherwise, and choose a \(K_{n,n}\)-free graph \(G\) of uncountable chromatic number whose vertex cardinality \(\kappa\) is minimal.

For \(A\in[\kappa]^n\), let

\[
N(A)=\{v\in\kappa:v\text{ is adjacent to every member of }A\}.
\]

Because \(G\) is \(K_{n,n}\)-free, \(|N(A)|<n\). Apply Lemma 1.4 to \(A\mapsto N(A)\), and let \(M_i,I_i\) be the resulting chain and layers.

By minimality of \(\kappa\), every induced graph \(G[I_i]\) is countably chromatic. Let \(G^\times\) consist of the edges joining distinct layers, and orient each such edge from the later endpoint toward the earlier endpoint. If \(v\in I_i\) had \(n\) earlier neighbours, those neighbours would form an \(n\)-set \(A\subseteq M_i\), and closure would give

\[
v\in N(A)\subseteq M_i,
\]

contrary to \(v\in I_i\). Thus the orientation of \(G^\times\) has outdegree at most \(n-1\). By Lemma 1.3, \(G^\times\) is finitely colourable.

Colour each \(G[I_i]\) with colours in \(\omega\), reusing the same palette on all layers, and combine this with a finite proper colouring of \(G^\times\). The resulting countable product colouring is proper on all of \(G\), a contradiction.
\end{proof}

\section{Bipartite expansions}\label{private-vertex-expansions-of-bipartite-graphs-are-obligatory}

\paragraph{Rainbow bipartite submatrices.}\label{rainbow-bipartite-submatrices}

\begin{lemma}[Rainbow bipartite submatrices]\label{lemma-3.1-rainbow-bipartite-submatrices}

For all positive integers \(n,t\), there exists \(q=q(n,t)\) with the following property. Suppose the edges of \(K_{q,q}\) are coloured, with an arbitrary colour set, so that at every vertex each individual colour occurs on fewer than \(t\) incident edges. Then there are \(n\)-element subsets \(X'\) and \(Y'\) of the two vertex classes such that all \(n^2\) edges between \(X'\) and \(Y'\) have distinct colours.

\end{lemma}
\begin{proof}
Choose independent uniformly random \(n\)-subsets \(X'\) and \(Y'\). For every unordered pair \(P\) of distinct same-coloured edges, let \(I_P\) be the indicator that both edges of \(P\) lie in the selected \(K_{n,n}\), and put
\[
 Z=\sum_P I_P.
\]
Thus \(Z=0\) exactly when all selected edges have distinct colours.

Let \(m_\gamma\) be the number of edges of colour \(\gamma\). The local bound gives

\[
m_\gamma\le(t-1)q.
\]

At a fixed vertex and for a fixed colour \(\gamma\), fewer than \(t\)
incident edges have colour \(\gamma\). Summing
\(\binom{d_{v,\gamma}}2\le(d_{v,\gamma})(t-1)/2\) over the \(2q\)
vertices shows that the number of same-coloured pairs sharing a vertex is
at most \((t-1)q^2\). The total number of same-coloured pairs is at most

\[
\sum_\gamma\binom{m_\gamma}{2}
\le \frac12\bigl(\max_\gamma m_\gamma\bigr)\sum_\gamma m_\gamma
\le \frac12(t-1)q^3.
\]

If the two edges share a vertex in the left class, their exact selection
probability is
\[
 \frac nq\frac{(n)_2}{(q)_2};
\]
the same formula holds when the shared vertex is on the right. If the two
edges are disjoint, the exact probability is
\[
 \left(\frac{(n)_2}{(q)_2}\right)^2.
\]
For \(q\ge2n\), these are at most \(2n^3/q^3\) and \(4n^4/q^4\),
respectively. By linearity of expectation,

\[
\mathbb{E}\!\left[Z\right]
 \le\frac{2(t-1)n^3}{q}+\frac{2(t-1)n^4}{q}.
\]

This is below \(1\) for sufficiently large \(q\). Since \(Z\) is
integer-valued and nonnegative, it cannot satisfy \(Z\ge1\) for every
choice when \(\mathbb{E}\!\left[Z\right]<1\). Hence some choice has \(Z=0\), and the
selected \(K_{n,n}\) is rainbow.
\end{proof}

\paragraph{The complete bipartite expansion atom.}\label{the-complete-bipartite-expansion-atom}
\begin{proposition}[The complete bipartite expansion atom]\label{proposition-3.2-the-complete-bipartite-expansion-atom}

For every positive integer \(n\), the triple system \(K_{n,n}^{+}\) is obligatory.

\end{proposition}
\begin{proof}
Suppose not. Choose a \(K_{n,n}^{+}\)-free triple system

\[
H=(\kappa,E)
\]

with \(\chi(H)>\aleph_0\) and with \(\kappa\) minimal.

For a pair \(xy\), write

\[
d_H(xy)=|\{z:\{x,y,z\}\in E\}|.
\]

Set \(t=3n^2+1\) and define a graph \(R\) on \(\kappa\) by

\[
xy\in E(R)\quad\Longleftrightarrow\quad d_H(xy)\ge t.
\]

\medskip\noindent\textbf{Claim 3.2.1 (the graph \(R\) is countably chromatic).}

Otherwise Lemma 2.1 gives a \(K_{n,n}\) in \(R\), with core vertex set \(C\), \(|C|=2n\). For each of its \(n^2\) graph edges \(xy\), choose a third vertex \(p_{xy}\) with

\[
\{x,y,p_{xy}\}\in E,
\]

avoiding \(C\) and all previously chosen private vertices. At every stage fewer than \(2n+n^2<t\) vertices are forbidden, while \(xy\) has at least \(t\) possible third vertices. Thus the \(p_{xy}\) can be chosen distinct and outside \(C\), producing a copy of \(K_{n,n}^{+}\), a contradiction. \(\square\)

Let

\[
E_1=\{e\in E:\text{some pair contained in }e\text{ is an edge of }R\},
\qquad E_2=E\setminus E_1.
\]

A proper colouring of \(R\) is a proper colouring of \((\kappa,E_1)\), since every \(E_1\)-edge contains an \(R\)-edge. Hence \(\chi(\kappa,E_1)\le\aleph_0\), and therefore

\[
\chi(\kappa,E_2)>\aleph_0. \tag{3.1}
\]

For a pair \(xy\), define

\[
\Phi(xy)=
\begin{cases}
\{z:\{x,y,z\}\in E\},&d_H(xy)<t,\\
\varnothing,&d_H(xy)\ge t.
\end{cases}
\]

Apply Lemma 1.4 to \(\Phi\), obtaining \(M_i,I_i\).

Every \(e\in E_2\) has at least two vertices in its highest layer. Indeed, suppose the highest layer met \(e\) only in \(z\in I_i\). Then the other pair \(xy=e\setminus\{z\}\) lies in \(M_i\). Since \(e\in E_2\), \(xy\notin R\), so \(d_H(xy)<t\). Closure gives \(z\in\Phi(xy)\subseteq M_i\), a contradiction.

Thus every \(E_2\)-edge is either contained in one layer, or has exactly two vertices in its highest layer \(I_i\) and its third vertex in \(M_i\). Let \(E_i^*\) be the \(E_2\)-edges contained in \(I_i\). Since \(|I_i|<\kappa\) and \(H\) is \(K_{n,n}^{+}\)-free, minimality of \(\kappa\) gives \(\chi(I_i,E_i^*)\le\aleph_0\). Consequently

\[
E^*=\bigcup_iE_i^*
\]

is countably chromatic. By (3.1), the remaining crossing-edge system

\[
E^{**}=E_2\setminus E^*
\]

has uncountable chromatic number.

For each \(i\), define a graph \(G_i\) on \(I_i\) by declaring \(xy\in E(G_i)\) when there is a \(\beta\in M_i\) such that \(\{x,y,\beta\}\in E^{**}\). If every \(G_i\) were countably chromatic, colour each \(I_i\) properly in \(G_i\), using the same countable palette on every layer. Every edge of \(E^{**}\) would then have differently coloured top-layer vertices, giving a countable colouring of \((\kappa,E^{**})\), a contradiction. Hence some \(G_i\) is uncountably chromatic.

Choose \(q=q(n,t)\) from Lemma 3.1. Lemma 2.1 gives a \(K_{q,q}\) in \(G_i\), with sides \(X,Y\). For every \(xy\in X\times Y\), choose

\[
\beta_{xy}\in M_i
\]

such that \(\{x,y,\beta_{xy}\}\in E^{**}\), and regard \(\beta_{xy}\) as the colour of the edge \(xy\).

This edge-colouring is locally \((t-1)\)-bounded. For example, if for fixed \(x\) and fixed \(\beta\) there were \(t\) different vertices \(y\) with \(\beta_{xy}=\beta\), then the pair \(\{x,\beta\}\) would have codegree at least \(t\). It would belong to \(R\), forcing the triples \(\{x,y,\beta\}\) into \(E_1\), contrary to their membership in \(E^{**}\subseteq E_2\). The same argument applies at vertices of \(Y\).

Lemma 3.1 gives \(n\)-sets \(X'\subseteq X\), \(Y'\subseteq Y\) for which the \(\beta_{xy}\), \(x\in X'\), \(y\in Y'\), are all distinct. They lie in \(M_i\), whereas \(X'\cup Y'\subseteq I_i\), so none is a core vertex. The triples

\[
\{x,y,\beta_{xy}\},
\qquad x\in X',\ y\in Y',
\]

form \(K_{n,n}^{+}\), the final contradiction.
\end{proof}

\begin{corollary}[Every finite bipartite expansion is obligatory]\label{corollary-3.3-every-finite-bipartite-expansion-is-obligatory}

If \(J\) is any finite bipartite graph, then \(J^+\) is obligatory.

\end{corollary}
\begin{proof}
Embed the two vertex classes of \(J\) in the two classes of some \(K_{n,n}\). Then \(J^+\) is a subhypergraph of \(K_{n,n}^+\). Obligatoriness is downward closed under non-induced containment: a copy of the larger finite system contains a copy of the smaller one. Apply Proposition 3.2.
\end{proof}

Corollary~3.3 also follows from \citet[Theorem~1.2]{reiher2024}.  The
argument above gives a direct proof in the notation used here.

\section{Closure properties}\label{closure-of-obligatoriness-under-the-operations-defining-mathcal-b}

Closure under disjoint unions and one-point amalgamations is known; see
\citet[pp.~233--238]{komjath2001} and the summary in
\citet[p.~1]{reiher2024}.  We include complete proofs because the subsequent
argument uses the rooted-abundance construction explicitly.

\paragraph{Disjoint-union closure.}\label{disjoint-union-closure}

\begin{lemma}[Disjoint-union closure]\label{lemma-4.1-disjoint-union-closure}

The class of finite obligatory triple systems is closed under finite disjoint unions.

\end{lemma}
\begin{proof}
Let \(F_1,F_2\) be obligatory and let \(H\) be uncountably chromatic. First find a copy of \(F_1\). Delete its finite vertex set. By Lemma 1.2 the remaining triple system is still uncountably chromatic, so it contains \(F_2\). The two copies are disjoint. Iteration proves the finite case.
\end{proof}

\paragraph{Rooted abundance and one-point amalgamation.}\label{rooted-abundance-and-one-point-amalgamation}

\begin{definition}[Rooted abundance]\label{definition-4.2-rooted-abundance}

Let \((F,r)\) be a finite rooted triple system with \(|V(F)|>1\). For a triple system \(H\) and \(m\ge1\), let \(R_m(F,r;H)\) be the set of vertices \(v\) for which there are \(m\) rooted copies of \(F\), all sending \(r\) to \(v\), whose off-root vertex sets are pairwise disjoint.

\end{definition}
\begin{lemma}[Rooted-abundance lemma]\label{lemma-4.3-rooted-abundance-lemma}

If \((F,r)\) is finite, \(|V(F)|>1\), and \(F\) is obligatory, and if
\(\chi(H)>\aleph_0\), then

\[
\chi\bigl(H[R_m(F,r;H)]\bigr)>\aleph_0.
\]

\end{lemma}
\begin{proof}
Put

\[
B=V(H)\setminus R_m(F,r;H).
\]

For each \(v\in B\), consider the family of off-root vertex sets of all
copies of \(F\) rooted at \(v\). Because
\(v\notin R_m(F,r;H)\), this family has matching number at most \(m-1\).
Choose a maximal pairwise disjoint subfamily. It contains at most \(m-1\)
members; let \(S_v\) be their union. Then

\[
|S_v|\le(m-1)(|V(F)|-1)=:D.
\]

By maximality, \(S_v\) meets the off-root part of every rooted copy at \(v\).
Otherwise that off-root set would be disjoint from every member of the
chosen subfamily and could be added to it.

Form a graph \(D_B\) on \(B\) by joining \(v\) to each member of \(S_v\cap B\). If an edge is generated in both directions, assign it to one of its generators, and orient it away from that generator. Every vertex has outdegree at most \(D\), so Lemma 1.3 gives a finite proper colouring of \(D_B\).

Let \(A\) be one colour class. Then

\[
S_v\cap A=\varnothing\qquad(v\in A).
\]

Consequently \(H[A]\) contains no copy of \(F\). Indeed, if \(\varphi\) were such a copy and \(v=\varphi(r)\), its off-root part would have to meet \(S_v\). All those off-root vertices lie in \(A\), contradicting \(S_v\cap A=\varnothing\). Since \(F\) is obligatory, every \(F\)-free triple system is countably chromatic. Thus each of the finitely many colour classes of \(D_B\) induces a countably chromatic subsystem. Tagging a vertex by its finite \(D_B\)-colour and its countable colour inside that class gives a countable proper colouring of \(H[B]\).

If \(H[R_m(F,r;H)]\) were countably chromatic as well, tagging the two sides would give a countable colouring of \(H\). This contradicts \(\chi(H)>\aleph_0\).
\end{proof}

\begin{proposition}[One-point amalgamation closure]\label{proposition-4.4-one-point-amalgamation-closure}

The class of finite obligatory triple systems is closed under one-point amalgamations.

\end{proposition}
\begin{proof}
Let \(F_1,F_2\) be obligatory, with selected vertices \(r_1,r_2\), and let \(F\) be their one-point amalgamation. If one factor consists only of its selected vertex, the assertion is immediate. Otherwise set

\[
m=|V(F_2)|.
\]

By Lemma 4.3, the set

\[
R=R_m(F_1,r_1;H)
\]

induces an uncountably chromatic subsystem of every uncountably chromatic \(H\). Hence \(H[R]\) contains a copy of \(F_2\); write \(\varphi\) for its embedding and let \(v=\varphi(r_2)\).

At \(v\) there are \(m\) rooted copies of \(F_1\) with pairwise disjoint off-root sets. The set \(\varphi(V(F_2)\setminus\{r_2\})\) has \(m-1\) vertices, so it meets at most \(m-1\) of those off-root sets. One rooted \(F_1\)-copy therefore avoids the \(F_2\)-copy outside \(v\). Their union is the required one-point amalgamation.
\end{proof}

Corollary 3.3, Lemma 4.1, Proposition 4.4, and the trivial edgeless case prove

\[
(2)\Longrightarrow(1). \tag{4.1}
\]

\section{The intrinsic decomposition}\label{equivalence-of-the-constructive-and-intrinsic-descriptions}

We state the selected-incidence decomposition and its running-intersection
argument in the explicit form used by the Lean development.  For a distinct
structural formulation, compare \citet[Sections~4--5]{li2026}.

We prove \((2)\Longleftrightarrow(3)\). Isolated vertices may be ignored by Lemma 1.1.

\paragraph{Forward preservation.}\label{forward-preservation}

\begin{proposition}[The intrinsic conditions are preserved by the generators and operations]\label{proposition-5.1-the-intrinsic-conditions-are-preserved-by-the-generators-and-operations}

Every generator of \(\mathcal B\) satisfies the intrinsic conditions, and the conditions are preserved under finite disjoint unions and one-point amalgamations.

\end{proposition}
\begin{proof}
For a bipartite graph \(J\), the expansion \(J^+\) is linear. Every hyperedge-node is joined to its private point by a bridge. A Berge cycle cannot use a private point, so the Berge cycles of \(J^+\) are precisely the ordinary graph cycles of \(J\), with the same lengths; these are even. Edgeless systems satisfy the conditions vacuously, and disjoint union plainly preserves them.

Consider a one-point amalgamation. Hyperedges belonging to different factors intersect only in the identified point, so linearity is preserved. Its Levi graph is the vertex-sum of the two Levi graphs at a point-node. Every simple cycle in a vertex-sum lies wholly in one factor: a cycle crossing into the other factor would have to pass through the cut vertex twice. Consequently bridges in the factors remain bridges, and every Berge cycle lies in one factor and retains its parity. This uses the standard cut-vertex structure of graph blocks \citep[Chapter~3]{diestel2017}; the corresponding hypergraph separation terminology is discussed in \citet{bahmanian2015}.
\end{proof}

Thus

\[
(2)\Longrightarrow(3). \tag{5.1}
\]

\paragraph{Bridge-block reconstruction.}\label{bridge-block-reconstruction}

\begin{proposition}[Bridge-block decomposition]\label{proposition-5.2-bridge-block-decomposition}

Suppose \(F\) is finite, has no isolated vertices, and satisfies the three intrinsic conditions. Then \(F\in\mathcal B\).

\end{proposition}
\begin{proof}
If \(E(F)=\varnothing\), then the assumption that \(F\) has no isolated vertices gives \(V(F)=\varnothing\), and \(F\) is an edgeless generator. Hence assume \(E(F)\ne\varnothing\). Put \(L=I(F)\), let \(B(L)\) be its set of bridges, and let \(\mathscr C\) be the set of connected components of \(L-B(L)\).

Consider a hyperedge-node \(e\). Since \(e\) is incident with a bridge and has degree three, it has at most two nonbridge incidences. It cannot have exactly one: a nonbridge incidence lies on a cycle, and that cycle must leave \(e\) through a second nonbridge incidence. Therefore every hyperedge-node has either no nonbridge incidences or exactly two.

\paragraph{The nontrivial components.}\label{the-nontrivial-components}

Let \(C\in\mathscr C\) contain a nonbridge incidence. Every hyperedge-node \(e\in C\) has exactly two point-neighbours in \(C\). Contract each such \(e\) to an ordinary graph edge between those two point-neighbours. This produces a finite graph \(J_C\).

The graph \(J_C\) is simple: parallel edges would correspond to two hyperedges of \(F\) sharing two points, contrary to linearity. Every cycle in \(J_C\) gives a Berge cycle of the same length in \(F\). Hence \(J_C\) has no odd cycle and is bipartite.

Each hyperedge-node \(e\in C\) has a third point-neighbour \(p_e\) across a bridge. These third points satisfy:

\begin{enumerate}
\def\labelenumi{\arabic{enumi}.}
\tightlist
\item
  \(p_e\notin C\); otherwise a path in \(C\) from \(p_e\) to \(e\), together with the incidence \(ep_e\), would place that incidence on a cycle.
\item
  If \(e\ne f\) lie in \(C\), then \(p_e\ne p_f\); otherwise a path in \(C\) from \(e\) to \(f\), together with \(ep_e\) and \(p_ff\), would place those bridge incidences on a cycle.
\end{enumerate}

Define a map from \(J_C^+\) to the subsystem \(P_C\) formed by these
hyperedges as follows. Send every graph vertex of \(J_C\) to the
corresponding point-node of \(C\), send the private vertex indexed by the
contracted edge associated with \(e\) to \(p_e\), and send that expanded edge
to the hyperedge \(e\) of \(F\). The two observations above show that the
vertex map is injective: the private points are pairwise distinct and lie
outside \(C\). Every \(e\) has exactly the two contracted endpoints and
\(p_e\) as its three vertices, so this map is incidence-preserving and
surjective on the vertices and hyperedges of \(P_C\). Hence it is an
isomorphism \(J_C^+\cong P_C\).

\paragraph{Components consisting of one hyperedge-node.}\label{components-consisting-of-one-hyperedge-node}

If a hyperedge-node \(e\) has no nonbridge incidence, it is an isolated component after the bridges are deleted. Its single triple is the expansion of a one-edge bipartite graph. Thus every hyperedge of \(F\) belongs to a piece \(P_C\) isomorphic to \(J^+\) for a finite bipartite graph \(J\). The hyperedge sets of these pieces partition \(E(F)\). Their vertex sets cover \(V(F)\): every point is incident with a hyperedge, and it either lies in that hyperedge-node's bridge-deleted component or is the point endpoint of a bridge leaving it.

\paragraph{Quotient forest.}\label{quotient-forest}

Contract every component of \(L-B(L)\). The original bridges become edges of a quotient multigraph \(T\). Two distinct bridges cannot join the same pair of bridge-deleted components: paths inside those two components, together with the two bridges, would form a cycle containing both, contradicting that they are bridges. Thus \(T\) is a simple graph.

To see that \(T\) is a forest, suppose that it contained a cycle. Within
each contracted component on this cycle, choose a path joining the
endpoints of the two incident quotient edges. Concatenating these paths
with the corresponding bridges gives a closed walk in \(L\) that traverses
each of those bridges exactly once. This is impossible: deleting a bridge
separates its endpoints, so every closed walk crosses the resulting cut an
even number of times. Hence \(T\) is a forest.

A vertex \(C\in V(T)=\mathscr C\) is called \emph{active} if the bridge-deleted component \(C\) contains a hyperedge-node. For an active \(C\), the piece \(P_C\) consists of the hyperedges represented in \(C\), all point-nodes of \(C\), and the point endpoints of bridge incidences leaving hyperedge-nodes in \(C\). An inactive component contains no hyperedge-node and hence no incidence of \(L-B(L)\), so it is a singleton point-node.

Every hyperedge now belongs to a bipartite expansion piece.  It remains to
order the pieces so that each new piece meets the assembled system in at most
one point.  The quotient forest records the attachment geometry needed for
this running-intersection property.

\begin{claim}[5.2.1. Running intersection]

Every component of \(T\) contains an active vertex, because \(F\) has no isolated vertices. Choose an active root in each component, and order the active vertices by nondecreasing depth, with arbitrary order among equal-depth vertices. When an active component \(C\) is added, its piece \(P_C\) meets the union of the earlier pieces in at most one point.

\end{claim}
\begin{proof}[Proof of the claim]
For a point \(p\in V(F)\), let \(X(p)\in\mathscr C\) be the component of \(L-B(L)\) containing the point-node \(p\). The active pieces containing \(p\) are precisely \(P_{X(p)}\), when \(X(p)\) is active, together with the pieces \(P_D\) at active neighbours \(D\) of \(X(p)\) whose corresponding original bridge has point endpoint \(p\). Thus the vertices of \(T\) indexing pieces that contain \(p\) are contained in the closed star centred at \(X(p)\).

Suppose that \(p\in P_C\) also lies in an earlier piece. If \(X(p)=C\),
every other piece containing \(p\) is indexed by a neighbour of \(C\). A tree
has only one neighbour of \(C\) at smaller depth, namely its parent, so only
that piece can be earlier. If \(X(p)\ne C\), then \(C\) is a neighbour of
\(X(p)\). The vertex \(X(p)\) cannot be a child of \(C\), because then both
\(X(p)\) and every other vertex in its closed star would be deeper than
\(C\). Hence \(X(p)\) is the parent of \(C\). In either case, \(p\) is the
point endpoint of the unique parent edge of \(C\). The same reasoning applies
to every point shared with an earlier piece, so all such points coincide with
this parent point \(p_C\), and

\[
P_C\cap\bigcup_{D\text{ earlier}}P_D\subseteq\{p_C\}.
\]

For a root active component the intersection is empty.
\end{proof}

Starting with the root piece in each component, add the active pieces in this
order. If the new intersection is empty, use a disjoint union. If it is the
singleton \(\{p_C\}\), use a one-point amalgamation at that point. Several
earlier pieces may already contain the same inactive attachment point, but
the new step still identifies only that one existing vertex.

The identity maps from the pieces into \(F\) assemble inductively to an
incidence isomorphism onto the union of the pieces already added.  By the
running-intersection claim, the next piece and the preceding union are either
disjoint or meet exactly at the point used in the amalgamation, so no
unintended identification is introduced.  The piece hyperedge sets are
disjoint and partition \(E(F)\), and their vertex sets cover \(V(F)\).
After all active pieces have been added, the resulting subsystem is \(F\).
Different components of \(T\) are combined by disjoint union, proving
\(F\in\mathcal B\).
\end{proof}

This proves

\[
(3)\Longrightarrow(2). \tag{5.2}
\]

Combining (5.1) and (5.2),

\[
(2)\Longleftrightarrow(3). \tag{5.3}
\]

\begin{theorem}[Canonical atom normal form]
\label{theorem-canonical-atom-normal-form}
Let $F$ be a finite triple system without isolated points satisfying the three
intrinsic conditions in Theorem~A.  Its hyperedges admit a canonical partition
into \emph{atoms} with the following properties.
\begin{enumerate}
\item Every atom is either one triple or $J^+$ for a finite $2$-connected
simple bipartite graph $J$.
\item Two distinct atoms meet in at most one point.  The bipartite incidence
graph between atoms and points belonging to at least two atoms is a forest.
\item The atom partition, the nontrivial cores $J$ up to graph isomorphism, and
the atom--point forest are determined by $F$.
\end{enumerate}
Conversely, every forest assembly of these atoms by disjoint unions and
one-point amalgamations satisfies the intrinsic conditions.
\end{theorem}
\begin{proof}
Let $L=I(F)$.  A \emph{cyclic block} is a maximal $2$-connected subgraph of
$L$ containing a cycle.  Since every hyperedge-node has degree three and is
incident with a bridge, its degree inside a cyclic block is two.  Suppressing
the hyperedge-nodes in such a block $C$ gives a finite graph $J_C$ on the
point-nodes of $C$.  It is simple by linearity and $2$-connected by standard
block theory \citep{whitney1932,diestel2017}.  Its cycles correspond to Berge cycles of the same lengths, so
$J_C$ is bipartite.

For a hyperedge-node $e\in C$, the third incidence is a bridge.  Its point
endpoint lies outside $C$, since otherwise that incidence would lie on a
cycle.  These third points are pairwise distinct as $e$ varies in $C$;
otherwise a path inside $C$ between two hyperedge-nodes, completed through the
common third point, would put both bridge incidences on a cycle.  The
hyperedges represented in $C$ therefore form exactly $J_C^+$.

A hyperedge-node outside every cyclic block has no nonbridge incidence and
contributes one single-triple atom.  Thus the atoms partition the hyperedges.
Two atoms cannot share two points, since paths inside the two atoms would
produce a Levi cycle crossing two cyclic blocks.  Similarly, a cycle in the
atom--shared-point incidence graph expands to a closed Levi walk and hence to
a Levi cycle crossing atom boundaries.  The incidence graph is therefore a
forest.  Cyclic blocks and the block--cut forest are canonical, proving the
forward statement.  The converse follows because an atomic expansion is
linear, has a private bridge at each hyperedge-node, and has only even Berge
cycles, while a forest of one-point sums preserves these properties.
\end{proof}

The graph primitives in this normal form are classical algorithmic objects:
bridges and bridge-connected components can be found by Tarjan's algorithm
\citep{tarjan1974}, while biconnected components admit a depth-first-search
construction \citep{tarjan1972}.  These algorithms are distinct from the
structural identification above; no verified running-time bound for the Lean
development is asserted.

Call a connected triple system with at least one hyperedge \emph{one-point
indecomposable} if it cannot be written as a one-point amalgamation of two
systems that both contain a hyperedge.

\begin{corollary}[Minimal generators and one-point indecomposables]
\label{corollary-minimal-generators-indecomposables}
Up to isomorphism, the class $\mathcal B$ is generated by finite edgeless
systems, one triple, and $J^+$ for finite $2$-connected simple bipartite graphs
$J$, under finite disjoint unions and one-point amalgamations.  A connected
reduced obligatory system with at least one hyperedge is
one-point indecomposable if and only if it is one triple or $J^+$ for a finite
$2$-connected simple bipartite graph $J$.
\end{corollary}
\begin{proof}
The theorem decomposes every reduced member of $\mathcal B$ into the displayed
atoms; isolated points are restored by an edgeless factor.  If the atom forest
has at least two atom-nodes, a leaf atom separates from the rest at its unique
shared point.  Conversely, one triple is indecomposable, and if $J$ is
$2$-connected then deleting any point of $J^+$ cannot separate its
hyperedge-nodes into two nonempty classes: a private point is a leaf, while a
core point leaves $J$ connected.  Hence $J^+$ is indecomposable.
\end{proof}

\section{The sequence lift}\label{the-sequence-lift-used-in-the-avoidance-constructions}

We use a one-apex construction specialised to $\omega_1$, presented through
base fibres and their support-incidence graph to match the formalisation.
\citet[Sections~3--4 and Theorem~4.6]{li2026} gives an alternative one-apex
and bridge-trace formulation.

\paragraph{The one-apex construction.}\label{the-one-apex-construction}

Let \(G=(X,A)\) be an arbitrary graph, where \(A=E(G)\) is also regarded as an alphabet. Let

\[
T=A^{<\omega_1}
\]

be the tree of all sequences of graph edges of countable ordinal length. For \(s\in T\) and \(x\in X\), write \(x_s=(s,x)\).

\begin{definition}[The one-apex lift]\label{definition-6.1-the-one-apex-lift}

Define the triple system \(\mathcal L(G)\) by

\[
V(\mathcal L(G))=T\times X.
\]

For \(s\in T\), an edge \(a=\{x,y\}\in A\), a sequence \(t\in T\) extending \(s^\frown a\), and \(z\in X\), put

\[
\{x_s,y_s,z_t\}\in E(\mathcal L(G)). \tag{6.1}
\]

The two vertices \(x_s,y_s\) form the base pair and \(z_t\) is the apex. Since \(t\) properly extends \(s\), these are three distinct vertices. In the resulting triple, the sequence node occurring twice uniquely identifies \(s\); its two graph labels identify the edge \(a=\{x,y\}\), and the remaining vertex identifies \(t,z\). Thus each hyperedge arises from unique data, up to interchanging \(x\) and \(y\), and the system is simple and \(3\)-uniform.

\end{definition}

\paragraph{Chromatic and finite-trace properties.}\label{chromatic-and-finite-trace-properties}

\begin{lemma}[Chromatic lower bound]\label{lemma-6.2-chromatic-lower-bound}

If \(\chi(G)>\aleph_0\), then

\[
\chi(\mathcal L(G))>\aleph_0.
\]

\end{lemma}
\begin{proof}
Suppose \(c:T\times X\to\omega\) is proper. For every \(s\in T\), the restriction \(x\mapsto c(x_s)\) is a countable colouring of \(G\). Hence there is an edge

\[
a_s=\{x_s^0,x_s^1\}\in E(G)
\]

whose two vertices at node \(s\) have a common colour \(k_s\).

Construct \(s_\alpha\in A^\alpha\) for \(\alpha<\omega_1\) by

\[
s_0=\varnothing,
\qquad s_{\alpha+1}=s_\alpha^\frown a_{s_\alpha},
\]

and at limit ordinals take unions.

If \(\alpha<\beta\), then \(s_\beta\) extends \(s_\alpha^\frown a_{s_\alpha}\). If \(k_{s_\alpha}=k_{s_\beta}\), take either endpoint \(z\) of the monochromatic edge at \(s_\beta\). By (6.1),

\[
\{(s_\alpha,x_{s_\alpha}^0),
  (s_\alpha,x_{s_\alpha}^1),
  (s_\beta,z)\}
\]

is a monochromatic hyperedge. Therefore the colours \(k_{s_\alpha}\), \(\alpha<\omega_1\), are pairwise distinct, impossible in \(\omega\).
\end{proof}

The following trace theorem uses the fibre decomposition described above.
Li gives a parallel bridge-trace formulation \citep[Theorem~4.6]{li2026}.

\begin{theorem}[Finite linear trace decomposition]\label{theorem-6.3-exact-finite-linear-trace-theorem}

Every finite, not necessarily induced, linear subhypergraph \(K\) of
\(\mathcal L(G)\) is obtainable, by finite disjoint unions and one-point
amalgamations, from systems \(J^+\) where \(J\) is a finite subgraph of
\(G\).

\end{theorem}
\begin{proof}
We use three local objects in the assembly.  A \emph{base node} is a sequence
\(s\) at which a lift hyperedge begins.  The two vertices with sequence
coordinate \(s\) are its \emph{core vertices}; the third vertex, whose
coordinate properly extends \(s\), is its \emph{apex}.  The \emph{base fibre}
\(K_s\) consists of all hyperedges of \(K\) based at \(s\), together with
their incident vertices.  A base node is \emph{active} when its fibre is
nonempty.  We first identify each fibre with an expansion piece and then show
that the fibres form a forest of one-point overlaps.

First set aside every isolated vertex of \(K\). Each is a one-vertex edgeless
factor, hence isomorphic to \(J^+\) for a one-vertex edgeless subgraph \(J\)
of \(G\), and may be restored by disjoint union at the end. We may therefore
assume that every vertex of \(K\) lies in a hyperedge.

Every hyperedge of \(\mathcal L(G)\) has a unique base node: exactly two of
its vertices have sequence coordinate \(s\), while the apex has a strictly
longer sequence coordinate. We group the hyperedges of \(K\) by this
coordinate and use the notation \(K_s\) from the dictionary.

\medskip
\noindent\textbf{Claim 6.3.1 (each base fibre is one expansion piece).}
For every active base node \(s\), there is a finite subgraph \(J_s\) of
\(G\) such that \(K_s\cong J_s^+\).

\smallskip
\begin{proof}[Proof of Claim 6.3.1]

For a fixed graph edge \(a=\{x,y\}\in E(G)\), at most one hyperedge of \(K_s\) can have base pair \(\{x_s,y_s\}\), since two such hyperedges would share two vertices and violate linearity. Let \(J_s\) be the finite subgraph of \(G\) whose edges are precisely the graph edges used by hyperedges of \(K_s\), with vertex set consisting of their endpoints.

Define a map \(J_s^+\to K_s\) by sending each graph vertex \(x\in V(J_s)\)
to the core vertex \(x_s\), and by sending the private vertex of
\(a=\{x,y\}\) to the apex of the unique hyperedge of \(K_s\) with base pair
\(\{x_s,y_s\}\). Distinct graph edges have distinct apices: if \(z_t\)
were the apex for \(a\) and \(a'\), the coordinate of \(t\) immediately
after \(s\) would have to equal both \(a\) and \(a'\). No apex is a core
vertex, because an apex has sequence coordinate strictly extending \(s\),
whereas every core vertex has coordinate \(s\). The map is therefore
injective, and by construction it is bijective on hyperedges and preserves
each incidence. Hence \(K_s\cong J_s^+\).
\end{proof}

\medskip
\noindent\textbf{Claim 6.3.2 (two fibres have at most one common point).}
For distinct active base nodes \(s\) and \(u\),
\[
|V(K_s)\cap V(K_u)|\le 1.
\]

\smallskip
\begin{proof}[Proof of Claim 6.3.2]
If a point \(z_t\) lies in both \(K_s\) and \(K_u\), then both \(s\) and
\(u\) are prefixes of \(t\), so they are comparable. Suppose
\(s\subsetneq u\). A common point cannot be a core point of \(K_s\); it is
an apex of \(K_s\). The graph edge used by its hyperedge in \(K_s\) is the
coordinate of \(u\) immediately following \(s\). Hence all points in
\(K_s\cap K_u\) would be apices of the same base edge of \(K_s\). There is
at most one such apex. Therefore

\[
|V(K_s)\cap V(K_u)|\le1. \tag{6.2}
\]
\end{proof}

Let \(Q\) be the bipartite incidence graph of the base fibres. Its left
class is the set of active base nodes. Its right class consists of points
that lie in at least two fibres. Join a base node \(s\) to a point \(p\)
exactly when \(p\in V(K_s)\).

\medskip
\noindent\textbf{Claim 6.3.3 (the overlap graph is a forest).}
The graph \(Q\) is acyclic.

\smallskip
\begin{proof}[Proof of Claim 6.3.3]
Suppose that \(Q\) has a cycle, and choose on it a base node \(s\) of
minimum sequence length. Its two neighbouring base nodes on the cycle are
proper descendants of \(s\). The two intervening shared points are distinct,
so they are distinct apices in \(K_s\), arising from two distinct graph edges
\(a\ne a'\). Thus the two neighbouring base nodes lie in different immediate
branches below \(s\), one beginning with \(a\) and the other with \(a'\).

Starting from one neighbour of \(s\), every base node on the remaining
part of the cycle lies in the same immediate branch below \(s\). Indeed,
consecutive base nodes share a point and are therefore comparable. For two
comparable proper descendants of \(s\), the first sequence coordinate after
\(s\) is the same, so a length-two step through a shared point cannot move
from one immediate branch to another. Induction along the remaining path
therefore keeps every base node in the first branch, contradicting that its
final node lies in the other branch. Hence \(Q\) is acyclic.
\end{proof}

Root each component of \(Q\) at a base node, and order its base nodes by
nondecreasing even distance from the root. For a nonroot base node \(s\), let
\(p_s\) be the point-node adjacent to \(s\) on the path to the root.

\begin{claim}[6.3.4. Base-fibre running intersection]

Every earlier factor that meets \(K_s\) meets it at \(p_s\).

\end{claim}
\begin{proof}[Proof of Claim 6.3.4]
Suppose \(K_u\) is earlier and contains a point
\(p\in K_s\). Then \(Q\) contains the length-two path \(s-p-u\). Let the
distance of \(s\) from the root be \(2d\). The point-node \(p\) has distance
either \(2d-1\) or \(2d+1\). In the second case \(u\) has distance \(2d+2\),
so it cannot be earlier. Thus \(p\) has distance \(2d-1\), and the uniqueness
of the path to the root in the tree \(Q\) gives \(p=p_s\). This also covers
an earlier sibling \(u\), which shares the same parent point \(p_s\).
\end{proof}

\medskip
\noindent\textbf{Claim 6.3.5 (the ordered fibres assemble to \(K\)).}
Let \(U_s\) be the union of the factors preceding \(K_s\). Equation (6.2)
and Claim 6.3.4 give
\[
V(K_s)\cap U_s\subseteq\{p_s\}.
\]
Starting with the root factor, add the remaining factors in this order by
disjoint unions or one-point amalgamations as appropriate. At each step the
natural maps into \(K\) identify exactly the displayed common point and
nothing else, so they assemble to an isomorphism onto the union of the
factors already added. Different components of \(Q\), and the isolated
factors removed at the start, are combined by disjoint union. This is the
running-intersection step: Claim 6.3.1 supplies the expansion pieces, while
Claims 6.3.2--6.3.4 certify that each new piece is attached along at most one
previous point.
\end{proof}

\begin{corollary}[Restrictions on finite linear traces]\label{corollary-6.4-restrictions-on-finite-linear-traces}

Every finite linear \(K\subseteq\mathcal L(G)\) satisfies:

\begin{enumerate}
\def\labelenumi{\arabic{enumi}.}
\tightlist
\item
  every hyperedge-node of \(I(K)\) is incident with a bridge;
\item
  every Berge cycle of \(K\) is contained in a single factor \(J_s^+\);
\item
  a Berge cycle of length \(\ell\) in \(K\) yields an ordinary cycle of length \(\ell\) in \(G\).
\end{enumerate}

\end{corollary}
\begin{proof}
In every \(J_s^+\) the incidence with the private vertex is a bridge. Bridgehood is preserved by one-point amalgamation. A simple Levi cycle cannot cross a one-point cut, so every Berge cycle lies in one factor. Berge cycles in \(J_s^+\) are exactly ordinary graph cycles in \(J_s\).
\end{proof}

\section{A classical high-odd-girth input}\label{uncountably-chromatic-graphs-of-large-odd-girth}

For the odd-cycle obstruction we use the following classical graph theorem
directly, without passing through a shift-graph surrogate.

\begin{theorem}[Erdős--Hajnal]\label{theorem-7.1-high-odd-girth}
For every positive integer $m$ there is a graph $G$ such that
\[
\chi(G)>\aleph_0
\]
and $G$ contains no odd cycle of length at most $m$.
\end{theorem}

\begin{proof}[Source]
Theorem~7.4 of \citet[p.~76]{erdoshajnal1966}, as quoted in
\citet[Theorem~C, p.~428]{egh1975}, gives for every positive integer $i$ an
uncountably chromatic graph containing no cycle $C_{2j+1}$ for $0<j<i$.
Choose $i$ so that $2i-1\ge m$.
\end{proof}

\section{Avoidance hosts}\label{the-avoidance-hosts}

The intrinsic criterion can fail in three ways: nonlinearity, a missing
bridge, or an odd Berge cycle.  \citet[Sections~2--4]{li2026} organises the
negative direction in a parallel way.  We prove the first two avoidance
constructions in our notation and apply the Erdős--Hajnal theorem from
Section~7 to the third.

We prove

\[
\neg(3)\Longrightarrow\neg(1).
\]

\paragraph{Nonlinearity obstruction.}\label{nonlinearity-obstruction}

\begin{proposition}[Avoidance of every nonlinear finite triple system]\label{proposition-8.1-avoidance-of-every-nonlinear-finite-triple-system}

There is an uncountably chromatic linear triple system. Consequently, every finite nonlinear triple system is non-obligatory.

\end{proposition}
\begin{proof}
We use the Erd\H{o}s--Rado relation
\citep[Theorem~4(i), formula~(95), p.~471]{erdosrado1956}

\[
(2^{\aleph_0})^+\longrightarrow(\aleph_1)^2_{\aleph_0}. \tag{8.1}
\]

Let \(\kappa=(2^{\aleph_0})^+\). Define a triple system \(T_\kappa\) whose vertices are the pairs in \([\kappa]^2\), and whose hyperedges are

\[
\bigl\{
\{\alpha,\beta\},
\{\alpha,\gamma\},
\{\beta,\gamma\}
\bigr\},
\qquad \alpha<\beta<\gamma<\kappa.
\]

This triple system is linear: two distinct triangles of a complete graph cannot share two graph edges, because two graph edges of a triangle determine the third edge and the three underlying vertices.

It is uncountably chromatic. A countable vertex-colouring of \(T_\kappa\) is a countable edge-colouring of \(K_\kappa\). By (8.1) there is a monochromatic subset of cardinality \(\aleph_1\), and in particular a monochromatic graph triangle. This is a monochromatic hyperedge of \(T_\kappa\).

If a finite triple system \(F\) is nonlinear, it has two distinct edges sharing at least two vertices. No injective embedding of those edges can exist in the linear host \(T_\kappa\). Thus \(T_\kappa\) is \(F\)-free.
\end{proof}

\paragraph{Missing-bridge obstruction.}\label{missing-bridge-obstruction}

\begin{proposition}[Avoidance when the bridge condition fails]\label{proposition-8.2-avoidance-when-the-bridge-condition-fails}

Let \(F\) be finite and linear, and suppose some hyperedge-node of \(I(F^\circ)\) is incident with no bridge. Then \(F\) is non-obligatory.

\end{proposition}
\begin{proof}
This is the $K_{\omega_1}$ specialisation of the sequence lift defined above.
Li gives an alternative missing-bridge obstruction \citep[Section~4]{li2026}.
By Lemma 1.1 it suffices to avoid \(F^\circ\),
so assume that \(F\) has no isolated vertices.
Take \(G=K_{\omega_1}\). Then \(\chi(G)>\aleph_0\), and Lemma 6.2 gives

\[
\chi(\mathcal L(G))>\aleph_0.
\]

If \(F\) appeared in \(\mathcal L(G)\), retain exactly the host hyperedges corresponding to the edges of \(F\). Because the embedding is injective and \(F\) is linear, these selected host hyperedges form a finite linear subhypergraph isomorphic to \(F\). Corollary 6.4 says that every hyperedge-node of such a trace is incident with a bridge, contradicting the chosen hyperedge of \(F\). Therefore \(\mathcal L(K_{\omega_1})\) is an uncountably chromatic \(F\)-free host.
\end{proof}

\paragraph{Odd-cycle obstruction.}\label{odd-cycle-obstruction}

\begin{proposition}[Avoidance of an odd Berge cycle]\label{proposition-8.3-avoidance-of-an-odd-berge-cycle}

Let \(F\) be finite and linear, and suppose \(F^\circ\) has an odd Berge cycle. Then \(F\) is non-obligatory.

\end{proposition}
\begin{proof}
Suppress isolated vertices and put $m=|E(F)|$.  By Theorem~7.1 choose a graph
$G$ with $\chi(G)>\aleph_0$ and no odd cycle of length at most $m$.  Lemma~6.2
gives $\chi(\mathcal L(G))>\aleph_0$.

Suppose that $F$ embeds in $\mathcal L(G)$ and retain exactly the host
hyperedges corresponding to the edges of $F$.  The selected trace is finite
and linear.  Corollary~6.4 localises every Berge cycle in one expansion fibre
and sends a Berge cycle of length $\ell$ to an ordinary cycle of the same
length in $G$.  The given odd Berge cycle has $\ell\le m$, contradicting the
choice of $G$.  Thus $\mathcal L(G)$ is an uncountably chromatic $F$-free
host.
\end{proof}

The three propositions cover every failure of the intrinsic conditions and
therefore prove

\[
\neg(3)\Longrightarrow\neg(1). \tag{8.2}
\]

\section{Proof of the classification}\label{conclusion}

Sections~3 and 4 establish \((2)\Longrightarrow(1)\), Section~5 establishes
\((2)\Longleftrightarrow(3)\), and Section~8 proves the contrapositive
\(\neg(3)\Longrightarrow\neg(1)\).  Together these implications prove all
equivalences in Theorem~A.  In particular,

\begingroup\small
\[
\begin{gathered}
F\text{ is obligatory}
\iff
F\in\mathcal B,\\[1mm]
F\in\mathcal B
\iff
\begin{array}{l}
F^\circ\text{ is linear},\\[1mm]
\text{every hyperedge-node of }I(F^\circ)
\text{ is incident with a bridge},\\[1mm]
\text{every Berge cycle of }F^\circ\text{ has even length}.
\end{array}
\end{gathered}
\]
\endgroup

All embeddings are non-induced, so additional host hyperedges on the image
vertices are irrelevant.  The proof uses only ZFC, graph-colouring
compactness \citep{debruijn1951}, the classical Erd\H{o}s--Hajnal high-odd-girth theorem \citep{erdoshajnal1966}, and the
standard Erd\H{o}s--Rado relation (8.1) \citep{erdosrado1956}.  It assumes neither CH nor GCH and
uses no forcing axiom or large-cardinal hypothesis.

\section{Finite parameter consequences}\label{finite-parameter-consequences}

In this section, a triple system is \emph{connected} when its Levi graph is
connected.  We first exclude isolated vertices.  For an integer $r\ge1$, set
\[
q(r)=\left\lceil2\sqrt r\right\rceil.
\]

\begin{proposition}[Edge-deletion form of the bridge condition]
\label{proposition-edge-deletion-bridge-condition}
Let $F$ be a finite triple system and let $e=\{x,y,z\}\in E(F)$.  The
hyperedge-node $e$ has no incident bridge in $I(F)$ if and only if $x,y,z$
lie in one connected component after the hyperedge $e$ is deleted and all
points are retained.  Consequently the bridge condition in Theorem~A is
equivalent to the following hypergraph-native condition:
\[
\begin{gathered}
\text{for every }e\in E(F^\circ),\text{ deleting }e\\
\text{separates its three points into at least two components.}
\end{gathered}
\]
\end{proposition}
\begin{proof}
Fix $x\in e$.  The Levi incidence $xe$ is not a bridge precisely when there
is an alternative path from $e$ to $x$ after that incidence is deleted.  A
simple such path leaves $e$ through $y$ or $z$ and thereafter lies entirely
in the Levi graph of $F-e$.  Thus $xe$ is non-bridging exactly when $x$ is
connected in $F-e$ to at least one of the other two points of $e$.

If all three incidences at $e$ are non-bridging, no one of $x,y,z$ can form a
singleton component among those three points.  A partition of three points
with this property has one block, so all three are connected in $F-e$.  The
converse follows by reversing the same paths.
\end{proof}

\begin{lemma}[Bipartite shadow]
\label{lemma-bipartite-shadow}
Let $F$ be an obligatory finite triple system with no isolated vertices and
at least one hyperedge.  There is a finite bipartite graph $J$, without
isolated vertices, such that
\[
|E(J)|=|E(F)|,\qquad |V(J)|=|V(F)|-|E(F)|,
\]
and $J$ has the same number of connected components as $F$.  Conversely, for
every finite bipartite graph $J$ without isolated vertices, the expansion
$J^+$ is obligatory and has these parameter relations.
\end{lemma}
\begin{proof}
By Theorem~A and Proposition~\ref{proposition-5.2-bridge-block-decomposition},
each connected component $H$ of $F$ is obtained by one-point amalgamations
from connected expansion pieces
\[
J_1^+,\ldots,J_k^+,
\]
where each $J_i$ is a connected bipartite graph.  Disconnected graph pieces
may simply be split into their connected components.  Write
$m_i=|E(J_i)|$ and $s_i=|V(J_i)|$.  Since $H$ is connected, its $k$ pieces
are joined by exactly $k-1$ one-point identifications.  Hence
\[
|E(H)|=\sum_{i=1}^k m_i,
\qquad
|V(H)|=\sum_{i=1}^k(m_i+s_i)-(k-1).
\tag{10.1}
\]

Take one-point sums of the ordinary graphs $J_1,\ldots,J_k$ along any tree on
the $k$ indices.  Before each identification, interchange the two
bipartition classes of the new factor if necessary.  The resulting graph
$J_H$ is connected and bipartite, and (10.1) gives
\[
|E(J_H)|=|E(H)|,
\qquad
|V(J_H)|=|V(H)|-|E(H)|.
\]
Taking the disjoint union of these shadows over the components of $F$ proves
the forward assertion.  The converse is the expansion-generator case of
Theorem~A, with the displayed counts immediate from the definition of $J^+$.
\end{proof}

\begin{theorem}[Exact order--size--component spectrum]
\label{theorem-order-size-component-spectrum}
Let $m\ge1$, $1\le c\le m$, and $n$ be integers.  There exists an obligatory
triple system $F$ with no isolated vertices, exactly $m$ hyperedges, exactly
$n$ vertices, and exactly $c$ connected components if and only if
\[
\boxed{
 m+2(c-1)+\left\lceil2\sqrt{m-c+1}\right\rceil
 \le n\le 2m+c.
}
\tag{10.2}
\]
Every integer in this interval occurs.
\end{theorem}
\begin{proof}
We begin with the corresponding elementary graph fact.  A connected simple
bipartite graph with $r\ge1$ edges and $s$ vertices exists if and only if
\[
q(r)\le s\le r+1.
\tag{10.3}
\]
Indeed, connectedness gives $r\ge s-1$.  If the bipartition has sizes $a,b$,
then
\[
r\le ab\le\left\lfloor\frac{s^2}{4}\right\rfloor,
\]
which is equivalent to $s\ge q(r)$.  Conversely,
$K_{\lfloor s/2\rfloor,\lceil s/2\rceil}$ contains a spanning tree, and
starting with that tree and adding arbitrary unused edges realises every edge
count between $s-1$ and $\lfloor s^2/4\rfloor$.

Apply Lemma~\ref{lemma-bipartite-shadow}.  Let the $c$ connected components of
the shadow have positive edge counts $m_1,\ldots,m_c$, where
$\sum_i m_i=m$.  If their orders are $s_i$, then (10.3) gives
\[
\sum_{i=1}^c q(m_i)\le \sum_{i=1}^c s_i\le m+c.
\tag{10.4}
\]
For positive integers $a,b$,
\[
q(a)+q(b)\ge q(a+b-1)+2.
\tag{10.5}
\]
To see this, note that
\[
\sqrt a+\sqrt b\ge\sqrt{a+b-1}+1,
\]
because, after squaring, the assertion reduces to
$(a-1)(b-1)\ge0$; taking ceilings gives (10.5).  Repeatedly
applying (10.5) to the left side of (10.4) yields
\[
\sum_{i=1}^c q(m_i)
\ge q(m-c+1)+2(c-1).
\]
The shadow has $n-m$ vertices, so these inequalities give (10.2).

For sharpness put $M=m-c+1$.  Take $c-1$ components equal to $K_2$.  By
(10.3), a final connected bipartite component with $M$ edges can have every
order from $q(M)$ through $M+1$.  Their disjoint union therefore realises
every shadow order from $q(M)+2(c-1)$ through $m+c$.  Expanding it adds
exactly $m$ private vertices and proves attainability of every integer in
(10.2).
\end{proof}

\begin{corollary}[Connected and unrestricted spectra]
\label{corollary-connected-order-size-spectrum}
For $m\ge1$:
\begin{enumerate}
\item a connected obligatory triple system without isolated vertices has
$m$ hyperedges and $n$ vertices if and only if
\[
m+\left\lceil2\sqrt m\right\rceil\le n\le2m+1;
\]
\item without prescribing the number of components, a reduced obligatory
triple system has $m$ hyperedges and $n$ vertices if and only if
\[
m+\left\lceil2\sqrt m\right\rceil\le n\le3m;
\]
\item if isolated vertices are permitted, an obligatory triple system with
$m$ hyperedges and $n$ vertices exists if and only if
\[
n\ge m+\left\lceil2\sqrt m\right\rceil.
\]
\end{enumerate}
\end{corollary}
\begin{proof}
The first assertion is Theorem~\ref{theorem-order-size-component-spectrum}
with $c=1$.  For the second, the connected interval covers through $2m+1$;
for $n>2m+1$, take $c=n-2m$ and use the upper endpoint $n=2m+c$ of (10.2).
The largest possible value is obtained from $m$ disjoint triples.  The final
assertion follows by adjoining arbitrary isolated vertices.
\end{proof}

\begin{corollary}[Size spectrum at fixed order]
\label{corollary-fixed-order-size-spectrum}
Let $c\ge1$ and $n\ge3c$.  A reduced obligatory triple system with exactly
$n$ vertices, $c$ connected components, and $m$ hyperedges exists if and only
if
\[
\boxed{
\left\lceil\frac{n-c}{2}\right\rceil
\le m\le
n-2c+4-\left\lceil2\sqrt{n-3c+4}\right\rceil.
}
\tag{10.6}
\]
Consequently, a reduced obligatory triple system with exactly $c$ connected
components and $n$ vertices exists if and only if
\[
n=3c,\qquad n=3c+2,\qquad\text{or}\qquad n\ge3c+4.
\tag{10.7}
\]
\end{corollary}
\begin{proof}
The upper-order inequality in (10.2) is equivalent to
$m\ge\lceil(n-c)/2\rceil$.  For the other inequality put
\[
M=m-c+1,\qquad N=n-3c+3.
\]
It becomes
\[
M+\left\lceil2\sqrt M\right\rceil\le N.
\tag{10.8}
\]
For integers $N\ge3$ and $M\ge1$, the largest $M$ satisfying (10.8) is
\[
U=N+2-\left\lceil2\sqrt{N+1}\right\rceil.
\tag{10.9}
\]
Indeed, write $k=\lceil2\sqrt{N+1}\rceil$.  Since
$k^2\ge4(N+1)$,
\[
4U=4N+8-4k\le(k-2)^2,
\]
so $\lceil2\sqrt U\rceil\le k-2$ and
$U+\lceil2\sqrt U\rceil\le N$.  On the other hand,
$(k-1)^2<4(N+1)$ gives
\[
(k-3)^2<4(N+3-k)=4(U+1),
\]
so $\lceil2\sqrt{U+1}\rceil\ge k-2$ and
$(U+1)+\lceil2\sqrt{U+1}\rceil\ge N+1$.  The left side of (10.8) is
strictly increasing in $M$, proving (10.9).  Translating back to $m$ gives
(10.6).

For (10.7), Corollary~\ref{corollary-connected-order-size-spectrum} shows that
a connected reduced obligatory system has order $3$, order $5$, or any order
at least $7$.  The first two statements come from $m=1,2$.  For $m\ge3$, the
connected intervals have no gaps because
$\lceil2\sqrt{m+1}\rceil\le m+1$ for $m+1\ge4$.  Thus every component has
order
\[
3+r,\qquad r\in\{0,2\}\cup[4,\infty).
\]
The sum of $c$ such increments is again $0$, $2$, or any integer at least
$4$: use one nonzero increment and take all remaining increments to be zero.
The increments $1$ and $3$ are impossible.  This proves (10.7).
\end{proof}

\begin{corollary}[Exact Levi cycle-rank spectrum]
\label{corollary-levi-cycle-rank-spectrum}
Let $F$ be reduced and obligatory, with $m$ hyperedges, $n$ vertices, and $c$
connected components.  The cyclomatic number of its Levi graph is
\[
\beta(I(F))=|E(I(F))|-|V(I(F))|+c=2m-n+c.
\tag{10.10}
\]
For fixed $m,c$, every integer in the interval
\[
0\le\beta(I(F))\le
m-c+2-\left\lceil2\sqrt{m-c+1}\right\rceil
\tag{10.11}
\]
occurs, and no other value occurs.  In particular,
\[
|V(F)|=2|E(F)|+c
\quad\Longleftrightarrow\quad
I(F)\text{ is a forest}.
\]
\end{corollary}
\begin{proof}
The Levi graph has $n+m$ vertices, $3m$ incidence edges, and $c$ components,
which gives (10.10).  Substitution of the two endpoints in (10.2) gives
(10.11), and every intermediate value occurs because every intermediate order
does.  A finite graph has cyclomatic number zero exactly when it is a forest.
\end{proof}

\begin{corollary}[Rigidity at balanced lower endpoints]
\label{corollary-balanced-endpoint-rigidity}
Let $F$ be connected, reduced, and obligatory.
\begin{enumerate}
\item If $|E(F)|=t^2$ and $|V(F)|=t^2+2t$, then
$F\cong K_{t,t}^+$.
\item If $|E(F)|=t(t+1)$ and $|V(F)|=t(t+1)+2t+1$, then
$F\cong K_{t,t+1}^+$.
\end{enumerate}
Here $t\ge1$, and the complete bipartite graphs are unique up to interchanging
the two classes.
\end{corollary}
\begin{proof}
In either case the shadow from Lemma~\ref{lemma-bipartite-shadow} attains
$|E(J)|=\lfloor|V(J)|^2/4\rfloor$.  Equality forces every possible cross edge
to be present and the two bipartition classes to be balanced, so the shadow is
$K_{t,t}$ or $K_{t,t+1}$, respectively.

For $t\ge2$ these complete bipartite graphs have no cut vertex.  On the other
hand, the shadow built from two or more connected expansion pieces is a
nontrivial one-point sum and has a cut vertex.  Hence the bridge-block
reconstruction has only one piece, and $F$ is the expansion of the displayed
complete bipartite graph.  The cases $t=1$ are immediate.
\end{proof}

\subsection{Canonical atoms and structural spectra}
\label{canonical-atoms-and-structural-spectra}

For a connected reduced obligatory system $F$, put
\[
s=|V(F)|-|E(F)|,
\qquad
\beta=2|E(F)|-|V(F)|+1.
\]
Thus $s$ is the order of a connected bipartite shadow and $\beta$ is its
cycle rank.  Let $a(F)$ denote the number of canonical atoms from
Theorem~\ref{theorem-canonical-atom-normal-form}.

\begin{lemma}[Minimum order at fixed cyclic rank]
\label{lemma-minimum-order-fixed-cyclic-rank}
Let $J$ be a finite $2$-connected simple bipartite graph of cycle rank
$r=|E(J)|-|V(J)|+1\ge1$.  Then
\[
|V(J)|\ge 2+q(r).
\tag{10.12}
\]
For $r\ge2$, every order $v\ge2+q(r)$ occurs.  For $r=1$, the possible orders
are exactly the even integers $v\ge4$.
\end{lemma}
\begin{proof}
Writing $v=|V(J)|$, the bipartite extremal bound gives
\[
r\le\left\lfloor\frac{v^2}{4}\right\rfloor-v+1
 =\left\lfloor\frac{(v-2)^2}{4}\right\rfloor,
\]
which is equivalent to (10.12).  If $v$ is even, start with the spanning cycle
$C_v$ in the balanced complete bipartite graph and add edges until the desired
rank is reached.  If $v$ is odd, start with $C_{v-1}$ and add the remaining
vertex with two neighbours in the opposite bipartition class; this is
$2$-connected and has rank two.  Add arbitrary missing cross-edges.  These
constructions realise every allowed $r\ge2$.  A $2$-connected graph of rank one
is a cycle, which is bipartite exactly when its order is even.
\end{proof}

\begin{theorem}[Exact canonical atom-count spectrum]
\label{theorem-exact-canonical-atom-count-spectrum}
Fix nonnegative integers $s,\beta$ satisfying either $\beta=0$ and $s\ge2$,
or $\beta\ge1$ and $s\ge2+q(\beta)$.  Let $F$ range over connected reduced
obligatory systems with at least one hyperedge, shadow order $s$ and cycle
rank $\beta$.  The possible values $k=a(F)$ are exactly
\[
\begin{array}{ll}
\beta=0:
  & k=s-1,\\[1mm]
\beta=1:
  & 1\le k\le s-3,\quad k\equiv s+1\pmod 2,\\[1mm]
\beta\ge2:
  & 1\le k\le s-1-q(\beta).
\end{array}
\tag{10.13}
\]
For $\beta\ge1$, the maximum is
\[
a_{\max}(s,\beta)=s-1-q(\beta).
\tag{10.14}
\]
Every maximiser has one minimum-order cyclic atom carrying the full rank
$\beta$, and every other atom is one triple.  The parity obstruction in the
rank-one line is the only gap in any atom-count spectrum.
\end{theorem}
\begin{proof}
Let the atom cores be $J_i$, with $v_i=|V(J_i)|$,
$e_i=|E(J_i)|$, and $r_i=e_i-v_i+1$.  Since a tree of $k$ atoms makes
$k-1$ one-point identifications,
\[
s=1+\sum_i(v_i-1),
\qquad
\beta=\sum_i r_i.
\tag{10.15}
\]
If $\beta=0$, every atom is one triple and (10.15) gives $k=s-1$.
If $\beta=1$, there is one cyclic atom, whose core is an even cycle of order
$v\ge4$; the remaining atoms are single triples, so
$s=k+v-1$.  This is exactly the middle line of (10.13).

Suppose $\beta\ge2$.  For positive integers $a,b$, the inequality
$\sqrt a+\sqrt b\ge\sqrt{a+b-1}+1$ follows after squaring from
$(a-1)(b-1)\ge0$.  Hence
$q(a)+q(b)\ge q(a+b-1)+2\ge q(a+b)+1$, where the second
inequality uses that $q$ increases by at most one between consecutive integers.
Together with Lemma~\ref{lemma-minimum-order-fixed-cyclic-rank}, repeated
application gives
\[
s\ge1+k+q(\beta),
\]
so $k\le s-1-q(\beta)$.  Conversely, for every integer in the displayed
interval choose one cyclic core of rank $\beta$ and order $s-k+1$, supplied by
the lemma, and attach $k-1$ single-triple atoms.  Equality in (10.14) forces
one positive-rank atom and minimum order, proving the rigidity statement.
\end{proof}

The rigidity in the maximum concerns the distribution of cyclic rank:
all positive rank lies in one minimum-order atom.  It does not assert that
the numerical parameters determine the graph-isomorphism type of that
atom's core.

For $s\ge3$, define
\[
\alpha(s)=
\begin{cases}
s,&s\text{ even},\\
s+1,&s\text{ odd},
\end{cases}
\qquad
\delta(s)=\left\lfloor\frac{(s-1)^2}{4}\right\rfloor+1.
\]

\begin{theorem}[Indecomposability phase diagram]
\label{theorem-indecomposability-phase-diagram}
Apart from the single triple $(|E(F)|,|V(F)|)=(1,3)$, a connected reduced
one-point-indecomposable obligatory system with shadow order $s$ and $m$
hyperedges exists if and only if
\[
s\ge4,
\qquad
\alpha(s)\le m\le\left\lfloor\frac{s^2}{4}\right\rfloor.
\tag{10.16}
\]
A connected reduced one-point-decomposable obligatory system exists if and
only if
\[
s\ge3,
\qquad
s-1\le m\le\delta(s).
\tag{10.17}
\]
Consequently, for $s\ge3$, the feasible connected parameters split into three exact zones:
\[
\begin{array}{rcl}
s-1\le m<\alpha(s)
 &\Longrightarrow& \text{every system is decomposable},\\
\alpha(s)\le m\le\delta(s)
 &\Longrightarrow& \text{both behaviours occur},\\
\delta(s)<m\le\lfloor s^2/4\rfloor
 &\Longrightarrow& \text{every system is indecomposable}.
\end{array}
\tag{10.18}
\]
For $m\ge2$, the last zone is equivalently
\[
s\le\left\lceil2\sqrt{m-1}\right\rceil.
\tag{10.19}
\]
\end{theorem}
\begin{proof}
By Corollary~\ref{corollary-minimal-generators-indecomposables}, the nontrivial
indecomposable systems are $J^+$ with $J$ finite, simple, bipartite and
$2$-connected.  If $|V(J)|=s$ and $|E(J)|=m$, minimum degree two gives
$m\ge s$; equality makes $J$ a cycle and therefore forces $s$ even.  The
upper bound is the usual $m\le\lfloor s^2/4\rfloor$.  Even cycles, and for odd
$s$ a cycle on $s-1$ vertices plus one new vertex with two neighbours, give
the lower endpoints; adding cross-edges fills the intervals.

A decomposable system admits a connected bipartite shadow with a cut vertex.
If its bipartition sizes are $a+b=s$ and the cut vertex lies in the $a$-side,
then deleting it gives components with sizes $(a_i,b_i)$ and
\[
m\le b+\sum_i a_i b_i
 \le b+(a-1)(b-1)=ab-a+1.
\]
The symmetric estimate gives $m\le ab-\min(a,b)+1\le\delta(s)$.  The bound is
sharp: take a connected bipartite graph on $s-1$ vertices with $m-1$ edges and
adjoin one leaf.  This also fills every value in (10.17).  The three zones and
the elementary inversion (10.19) follow.
\end{proof}

\begin{corollary}[Rigidity on the structural boundaries]
\label{corollary-structural-boundary-rigidity}
Let $F$ be connected, reduced and obligatory, with shadow order $s$ and
$m=|E(F)|$ hyperedges.
\begin{enumerate}
\item If $F$ is indecomposable, $s$ is even and $m=s$, then
$F\cong C_s^+$.
\item If $F$ is indecomposable, $s$ is odd and $m=s+1$, then
$F\cong\Theta(2a,2b,2c)^+$ for positive $a\le b\le c$ with
$a+b+c=(s+1)/2$.
\item If $F$ is indecomposable and $m=\lfloor s^2/4\rfloor$, then
$F\cong K_{\lfloor s/2\rfloor,\lceil s/2\rceil}^+$.
\item If $F$ is decomposable and $m=\delta(s)$, then $F$ is a one-point
amalgamation of $K_{a,b}^+$ and one triple, where
$a+b=s-1$ and $|a-b|\le1$.
\end{enumerate}
Here $\Theta(p,q,r)$ is the theta graph with three internally disjoint paths
of lengths $p,q,r$ and common endpoints.
\end{corollary}
\begin{proof}
The first and third assertions are the equality cases of minimum degree and
the bipartite extremal bound.  In the second case
$\sum_v(\deg v-2)=2$.  A single degree-four vertex would be a cut vertex, so
there are exactly two degree-three vertices and all others have degree two;
the three resulting paths have the same parity.  Their total length is
$s+1$, which is even, so all three paths are even.  For the fourth assertion
when $s\ge5$, choose the shadow by assembling the canonical atom cores.  Its
cyclic blocks are precisely the nontrivial canonical cores.  Equality in the cut-vertex estimate
forces a balanced complete bipartite graph on $s-1$ vertices plus one leaf,
and hence the displayed complete-bipartite atom and one singleton atom.
The tree cases $s=3,4$ are immediate.
\end{proof}

\begin{proposition}[Componentwise atom-count spectrum]
\label{proposition-componentwise-atom-count-spectrum}
Let $F$ be reduced and obligatory with at least one hyperedge and $c\ge1$
connected components, shadow order
$s=|V(F)|-|E(F)|$, cycle rank
$\beta=2|E(F)|-|V(F)|+c$, and $k$ canonical atoms.  For every feasible
$(s,\beta,c)$, the possible values of $k$ are exactly
\[
\begin{array}{ll}
\beta=0:
  & k=s-c,\\[1mm]
\beta=1:
  & c\le k\le s-c-2,\quad k\equiv s-c\pmod2,\\[1mm]
\beta\ge2:
  & c\le k\le s-c-q(\beta).
\end{array}
\tag{10.20}
\]
The parameters themselves are feasible exactly when $s\ge2c$ for $\beta=0$
and $s\ge2c+q(\beta)$ for $\beta\ge1$.
\end{proposition}
\begin{proof}
If the atom cores are $J_i$, a forest with $k$ atom-nodes and $c$ components
makes $k-c$ identifications.  Hence
\[
s=c+\sum_i(|V(J_i)|-1),
\qquad
\beta=\sum_i\bigl(|E(J_i)|-|V(J_i)|+1\bigr).
\]
The connected proof applies componentwise.  Necessity gives the lower atom
floor $k\ge c$ and the three upper formulas.  For sufficiency, take $c-1$
single-triple components and realise the corresponding connected atom count
in the remaining component.  The feasibility thresholds follow by
concentrating all positive rank in one minimum-order cyclic atom.
\end{proof}

\begin{remark}[Blocks and the earlier reduction]
\label{remark-komjath-block-reduction}
Komj\'ath proved that a finite triple system is obligatory if and only if all
of its $2$-connected components are obligatory \citep{komjath2001}.  General
hypergraph block decomposition and its relation to the incidence graph are
classical; see \citet{whitney1932}, \citet[Chapter~3]{diestel2017}, and, for the hypergraph formulation, \citet{bahmanian2015}.  For the reduced
$3$-uniform systems considered here, the canonical atoms in
Theorem~\ref{theorem-canonical-atom-normal-form} are precisely these maximal
nonseparable edge-supported blocks.  Thus the new content of the canonical
atom theorem is not the existence of a block decomposition, but the exact
identification of the block types that can occur in an obligatory triple
system, together with the quantitative spectra derived from them.
\end{remark}

\begin{corollary}[Block formulation of the classification]
\label{corollary-block-formulation}
Let $F$ be a finite triple system.  After isolated points are deleted, $F$ is
obligatory if and only if every nonempty block of $F^\circ$ is either one
triple or $J^+$ for a finite $2$-connected simple bipartite graph $J$.
\end{corollary}
\begin{proof}
If $E(F^\circ)=\varnothing$, then $F^\circ$ is empty and the assertion follows
from Lemma~\ref{lemma-1.1-isolated-vertex-reduction}.  Otherwise suppose first
that $F$ is obligatory.  Lemma~\ref{lemma-1.1-isolated-vertex-reduction}
gives that $F^\circ$ is obligatory, and Komj\'ath's reduction gives
obligatoriness of each of its blocks.  A block is one-point indecomposable,
so Corollary~\ref{corollary-minimal-generators-indecomposables} gives exactly
the displayed list.

Conversely, each displayed block is obligatory by the positive expansion
result (the one-triple case included).  Komj\'ath's reduction therefore makes
$F^\circ$ obligatory, and Lemma~\ref{lemma-1.1-isolated-vertex-reduction}
restores the isolated points.
\end{proof}

\begin{lemma}[Local product for one-point decompositions]
\label{lemma-local-decomposition-lattice-product}
Let $F$ be reduced and obligatory with at least one hyperedge.  Let
$\mathcal A(F)$ be its canonical atoms, let $S(F)$ be the points belonging to
at least two atoms, and put
\[
 \mu(p)=|\{A\in\mathcal A(F):p\in V(A)\}|.
\]
Let $\mathcal D(F)$ be the set of partitions of $E(F)$ into nonempty connected
supported pieces such that two pieces meet in at most one point and the
piece--shared-point incidence graph is a forest, ordered by refinement.  Then
$\mathcal D(F)$ is a lattice and
\[
 \mathcal D(F)\cong\prod_{p\in S(F)}\Pi_{\mu(p)},
 \tag{\(\dagger\)}
\]
where $\Pi_r$ denotes the partition lattice of an $r$-element set.  If $F$
has $k$ canonical atoms and $c$ connected components, then
\[
 \operatorname{rk}\mathcal D(F)=k-c
   =\sum_{p\in S(F)}(\mu(p)-1).
 \tag{\(\ddagger\)}
\]
The empty product in \((\dagger)\) is the one-element lattice.
\end{lemma}
\begin{proof}
A supported one-point forest decomposition cannot split a block: if two
hyperedges of one block lay in different pieces, the unique path between
those pieces in the piece--shared-point forest would contain a point whose
deletion separates the two hyperedges, contradicting nonseparability of the
block.  Thus every element of $\mathcal D(F)$ is a coarsening of the canonical
atom decomposition.

For a shared point $p$, restrict such a coarsening to the set of the
$\mu(p)$ atoms through $p$; this gives an element of $\Pi_{\mu(p)}$.
Conversely, choose independently a partition of the atoms through every
shared point and take the equivalence relation on all atoms generated by
these local identifications.  To track its classes without contracting
overlapping subtrees, replace the star at each shared point $p$ by the
following two-level tree.  Retain a central node $p$, insert one node $(p,B)$
for each block $B$ of the prescribed local partition, join $p$ to $(p,B)$,
and join $(p,B)$ to each atom in $B$.  Replacing a star of a forest by a tree
with the same attachment nodes preserves acyclicity, so the resulting
auxiliary graph is a forest.

Delete the edges from the central nodes $p$ to their nodes $(p,B)$.  The
components containing atoms are vertex-disjoint and correspond
exactly to the generated equivalence classes; their atom unions are
connected supported pieces.  Two different blocks at one $p$ cannot lie in
the same component, since a path between them, together with their two edges
to $p$, would give a cycle in the auxiliary graph.  Contract these
components in the full auxiliary graph, with the deleted edges restored,
while retaining the central nodes $p$.  There are
therefore no repeated incidences between a piece and $p$.  Delete each
degree-one central node, which is now internal to one piece.  The resulting
graph is precisely the piece--shared-point incidence graph and is still a
forest.  In particular, two pieces cannot meet at two shared points.
Restriction recovers the prescribed local partitions.  Conversely, a
connected supported piece is recovered from its local identifications,
since its atom nodes are connected through the shared points it contains.
Thus the constructions are inverse and preserve refinement, proving
\((\dagger)\).  For the related theory of connected set partitions and bond
lattices, see \citet{simon2011} and \citet[Lecture~2]{stanley2007}.
The factors here are full partition lattices on the atoms through each
shared point, not the bond lattice of the incidence forest itself.
Binary joints contribute two-element-chain factors, whereas three atoms
through one point already give the diamond $\Pi_3\cong M_3$, the familiar
obstruction to distributivity \citep[Chapter~II, Theorem~101]{gratzer2011}.

Finally the canonical incidence forest has $k+|S(F)|$ vertices, $c$
components and $\sum_p\mu(p)$ edges.  Hence
\[
 \sum_p(\mu(p)-1)=k-c,
\]
which is also the rank of the product lattice by the standard rank formula
for partition lattices and their products \citep[Chapter~3]{stanley2012}.
\end{proof}

\begin{lemma}[Capacity-safe realization of every attachment profile]
\label{lemma-capacity-safe-profile-realization}
Fix pairwise disjoint canonical atoms $A_1,\ldots,A_k$ and an integer
$1\le c\le k$.  For every integer partition
\[
 \lambda=(\lambda_1,\ldots,\lambda_t)\vdash k-c
\]
(with the empty partition allowed when $k=c$), there is a forest assembly of
the $A_i$ into exactly $c$ connected components whose shared-point excess
profile is $\lambda$; that is, its nontrivial shared points have
multiplicities $\lambda_i+1$.  Every atom is incident with at most two
distinct shared points.  Consequently, for a fixed atom list and fixed $c$,
all these assemblies have the same order, size, shadow order and cycle rank.
\end{lemma}
\begin{proof}
Put $N=k-c$.  If $N=0$, take the $k=c$ atoms as separate components.  Assume
$N>0$.  Use $N+1$ atoms in one nontrivial component and leave the remaining
$c-1$ atoms as singleton components.

If $t=1$, identify one chosen vertex from each of the $N+1=\lambda_1+1$
active atoms to one common shared point.  If $t\ge2$, arrange shared points
$p_1,\ldots,p_t$ in a chain.  Use $t-1$ atoms as connectors, the $i$th
connector meeting $p_i$ and $p_{i+1}$.  The two endpoint shared points already
have one incident connector and the internal shared points have two.  Supply
the remaining incidences with leaf atoms so that
$\deg(p_i)=\lambda_i+1$.  The number of required leaf atoms is
\[
 \lambda_1+\lambda_t+\sum_{i=2}^{t-1}(\lambda_i-1)
   =N-t+2,
\]
so together with the $t-1$ connectors this uses exactly $N+1$ atoms.  Every
connector atom uses two distinct vertices and every leaf atom one; canonical
atoms have at least three vertices, so the identifications can be made with
distinct ports inside each atom.  The atom--shared-point incidence graph is a
tree and has the required multiplicities.

Disjoint union and one-point amalgamation preserve obligatoriness by
Lemma~\ref{lemma-4.1-disjoint-union-closure} and
Proposition~\ref{proposition-4.4-one-point-amalgamation-closure}.
Each shared point separates the incident atom branches in this forest
assembly.  A nonseparable supported block therefore cannot contain edges
from two different atoms: the intervening shared point would separate
those edges.  Each supplied atom is itself one-point indecomposable, so its
image is a maximal nonseparable supported block.  By
Theorem~\ref{theorem-canonical-atom-normal-form}, these images are precisely
the canonical atoms of the assembled system.  Thus the constructed profile
is the canonical attachment profile, not merely an auxiliary incidence
pattern.

Any forest assembly of the fixed $k$ atoms into $c$ components makes exactly
$k-c$ vertex identifications.  Thus
\[
 |E(F)|=\sum_i|E(A_i)|,
 \qquad
 |V(F)|=\sum_i|V(A_i)|-(k-c),
\]
so the stated global parameters do not depend on the chosen profile.
\end{proof}

\begin{corollary}[Exact one-point factorization-lattice spectrum]
\label{corollary-factorization-lattice-spectrum}
Fix a feasible triple $(s,\beta,c)$ for reduced obligatory systems with at
least one hyperedge, where
\[
 s=|V(F)|-|E(F)|,
 \qquad
 \beta=2|E(F)|-|V(F)|+c,
 \qquad
 q(r)=\left\lceil2\sqrt r\right\rceil.
\]
As $F$ ranges over all such systems with these fixed parameters, put $k$ for
its number of canonical atoms and $N=k-c$.  The possible values of $N$ are
exactly
\[
\begin{array}{ll}
\beta=0:&N=s-2c,\\[1mm]
\beta=1:&0\le N\le s-2c-2,\quad N\equiv s-2c\pmod2,\\[1mm]
\beta\ge2:&0\le N\le s-2c-q(\beta).
\end{array}
\tag{\(\star\)}
\]
For every allowed $N$, the possible lattice isomorphism types of
$\mathcal D(F)$ are precisely
\[
 \prod_{i=1}^t\Pi_{\lambda_i+1},
 \qquad
 \lambda=(\lambda_1,\ldots,\lambda_t)\vdash N,
 \tag{\(\star\star\)}
\]
where for $N=0$ the empty partition gives the one-element lattice.  Every
lattice in \((\star\star)\) occurs, and different partitions $\lambda$ give
nonisomorphic lattices.
\end{corollary}
\begin{proof}
The rank spectrum \((\star)\) is
Proposition~\ref{proposition-componentwise-atom-count-spectrum} after
subtracting $c$ from the atom count $k$.  Fix an allowed $N$ and choose one
system with $k=N+c$ atoms and the prescribed global parameters.  By
Lemma~\ref{lemma-capacity-safe-profile-realization}, its fixed atom list can be
reassembled with the same $c$ and the same global parameters to realize every
partition $\lambda\vdash N$.  Lemma~\ref{lemma-local-decomposition-lattice-product}
then gives exactly the product in \((\star\star)\).  Conversely every system
has an excess profile partitioning $N$, so there are no other types.

Finally, the standard characteristic polynomial of the partition lattice
\citep[Chapter~3]{stanley2012} is
\[
 \chi_{\Pi_{r}}(x)=\prod_{j=1}^{r-1}(x-j),
\]
and characteristic polynomials multiply under direct products, by the
product rule for M\"obius functions \citep{rota1964,stanley2012}.
For the equivalent braid-arrangement viewpoint, see
\citet[Lectures~1--2]{stanley2007}.  Hence the
multiplicity of the root $j$ in the characteristic polynomial of the lattice
in \((\star\star)\) is
\[
 |\{i:\lambda_i\ge j\}|.
\]
These multiplicities recover the partition $\lambda$, so distinct partitions
give nonisomorphic lattices.
\end{proof}

Uniqueness of indecomposable direct factors also has a classical incidence-poset
formulation; compare \citet[Lemma~6.1]{schmitt1994}.  The argument here recovers
the attachment profile directly from characteristic-polynomial root multiplicities.

The numerical and lattice spectra record different parts of the same
decomposition.  The shadow order, cyclic rank and component count restrict
the possible excess $N=k-c$.  For each allowed $N$, the attachment profile
$\lambda\vdash N$ determines the product of partition lattices, and every
such profile is realizable.  These assertions classify the available
factorization-lattice types; they do not identify all systems with the same
numerical parameters.

\subsection{Cyclic costs and assembly freedom}
\label{cyclic-costs-and-assembly-freedom}

The atom-count bound and the local product measure complementary constraints.
The first records how much cyclic structure costs; the second records the
decomposition choices left after the atoms have been fixed.

\begin{proposition}[Exact cyclic-cost budget]
\label{proposition-cyclic-cost-budget}
Let $F$ be connected, reduced and obligatory, with positive Levi cycle rank
$\beta$.  Put $s=|V(F)|-|E(F)|$, let $k$ be its canonical atom count, and
set $N=k-1$.  For the positive-rank atoms, write $r_i$ for the core rank,
$v_i$ for the core order, and $h$ for their number.  Define
\[
 C=\sum_{r_i>0}q(r_i)-q(\beta),
 \qquad
 L=\sum_{r_i>0}\bigl(v_i-2-q(r_i)\bigr).
\]
Then
\[
 N+C+L=s-2-q(\beta),
 \qquad C\ge h-1,\quad L\ge0.
\]
In particular,
\[
 N+(h-1)+L\le s-2-q(\beta).
\]
The cost $C+L$ is zero exactly when the atom count is maximal: there is one
minimum-order cyclic atom carrying rank $\beta$, and all other atoms are
single triples.
\end{proposition}
\begin{proof}
Zero-rank atoms have core order two.  Equation (10.15) therefore gives
\[
 s=1+k+\sum_{r_i>0}q(r_i)+L,
 \qquad \sum_{r_i>0}r_i=\beta.
\]
Substituting the definitions of $N$ and $C$ gives the identity.
Lemma~\ref{lemma-minimum-order-fixed-cyclic-rank} gives $L\ge0$.
Since $\beta>0$, there is at least one positive-rank atom.  The strict
aggregation inequality in the proof of
Theorem~\ref{theorem-exact-canonical-atom-count-spectrum} gives $C\ge h-1$.
Thus zero cost forces $h=1$ and zero core-order slack; the converse follows
from the same identity.
\end{proof}

Since $N=\operatorname{rk}\mathcal D(F)$, the budget also describes the rank
available to the factorization lattice.  Splitting cyclic rank among atoms
and increasing their core orders reduce this rank.  It is a numerical
constraint, not a graph-edit-distance stability theorem or a uniqueness
statement about the cores.

We use the usual second-kind Stirling numbers: they count partitions into
nonempty blocks, with $S(0,0)=1$ \citep[Section~5.1]{comtet1974}.

\begin{corollary}[Sharp counts at a prescribed piece number]
\label{corollary-sharp-piece-counts}
Let $F$ be reduced and obligatory with at least one hyperedge, $k$ canonical
atoms and $c$ connected components.  Put $N=k-c$ and let
$\lambda=(\lambda_1,\ldots,\lambda_t)\vdash N$ be its attachment excess
profile.  For $0\le j\le N$, define
\[
 d_j(F)=|\{D\in\mathcal D(F):|D|=k-j\}|,
\]
where $|D|$ counts the actual nonempty original-edge pieces.  Write $S(a,b)$
for the number of partitions of an $a$-element set into $b$ blocks.  Then
\[
 \sum_{j=0}^{N}d_j(F)z^j
 =\prod_{i=1}^{t}
   \left(\sum_{r=0}^{\lambda_i}
     S(\lambda_i+1,\lambda_i+1-r)z^r\right),
\]
and
\[
 \binom Nj\le d_j(F)\le S(N+1,N+1-j)
 \qquad(0\le j\le N).
\]
If $N\ge2$, then for every $1\le j\le N-1$ the lower equality holds exactly
when all shared points are binary, that is, $\lambda=(1,\ldots,1)$; the
upper equality holds exactly when $\lambda=(N)$.  Each bound is attained
at every $j$ by its corresponding extremal profile, even with the atom list
and component count fixed.  For $N=0$ the generating polynomial is $1$, and for $N=1$ the
two extremal profiles coincide.
\end{corollary}
\begin{proof}
Let $\pi_i$ be the local partition at the $i$th shared point.  In the
auxiliary forest from
Lemma~\ref{lemma-local-decomposition-lattice-product}, delete the central
point-nodes and their incident edges.  The remaining forest has
$k+\sum_i|\pi_i|$ vertices and $\sum_i(\lambda_i+1)$ edges.  Every component
contains an atom and corresponds to one actual piece.  Hence
\[
 |D|=k-\sum_i\bigl(\lambda_i+1-|\pi_i|\bigr).
\]
Independent local choices give the generating polynomial.  This is the
standard partition-lattice enumeration
\citep[Example~3.10.4]{stanley2012}; connected-partition polynomials and their
articulation factorization are discussed in \citet[Sections~3--5]{simon2011}.

For the lower bound, give the $i$th local set one anchor and $\lambda_i$
ordinary elements.  A $j$-element subset of the total $N$ ordinary elements
defines local partitions by joining its elements to their respective
anchors and leaving all others singleton.  This is an injection into the
choices counted by $d_j(F)$.

For the upper bound, identify all local anchors to one anchor on a set of
$N+1$ elements.  Fuse the anchor blocks of the local partitions and retain
every other block.  Restriction to each anchored local set recovers its
partition, so the map is injective.  A tuple of total excess $j$ gives
$N+1-j$ blocks, proving the upper bound.  This construction also shows that
merging two positive profile parts never decreases any coefficient.

For strictness in the lower bound, suppose some $\lambda_i\ge2$.  Form an
anchor-free pair in that group and choose $j-1$ of the remaining $N-2$
ordinary elements to join their local anchors.  This tuple has excess $j$
but is outside the lower image.  For strictness in the upper bound, suppose
there are at least two profile groups.  Choose $j+1$ ordinary elements
meeting two groups, put them in one anchor-free block and leave all other
elements singleton.  This partition has $N+1-j$ blocks but is outside the
upper image.  Both choices are possible for $1\le j\le N-1$.
The two extremal profiles attain the respective bounds.  Finally,
Lemma~\ref{lemma-capacity-safe-profile-realization} realizes both profiles
using the same atoms and component count, with all global parameters
preserved.
\end{proof}

For example, take $C_4^+$ and two single triples.  A chain assembly with two
binary joints has decomposition polynomial $1+2z+z^2$; an assembly with both
triples attached at one point of $C_4^+$ has polynomial $1+3z+z^2$.
Both systems have twelve vertices, six hyperedges, three canonical atoms
and cycle rank one, and both maximize the atom count for these parameters.
Their four and five decomposition choices differ only because of attachment
geometry.  Thus extremal cyclic concentration does not eliminate assembly
freedom.  The counting inequalities above are elementary partition-lattice
consequences; their role here is to quantify that freedom, not to assert a
new general enumeration theorem.
The individual factors are reversed rows of second-kind Stirling numbers,
whose generating-function and concavity properties are classical
\citep{lieb1968}.  The bounds above instead compare different attachment
profiles at fixed total excess.
Pitman's probabilistic representations and coefficient bounds give further
context for these classical Stirling rows \citep{pitman1997}.  Our question
is which products of the rows are realized by attachment profiles, not a
new general identity for Stirling numbers.

\section{Quantitative comparisons for graph cores}
\label{sec:quantitative-cores}

The results in this section concern the graph $J$ underlying an expansion
$J^+$, not uncountably chromatic hypergraph hosts.  If $J$ is a simple
bipartite subdivision of $K_4$, its six path lengths $L_{ij}\ge1$ satisfy
$L_{ij}\equiv\varepsilon_i+\varepsilon_j\pmod2$ for some branch colours
$\varepsilon_i\in\{0,1\}$.  Its path interiors are mutually disjoint and
$|V(J)|=|E(J)|-2$.  The canonical-atom theorem admits every such $J^+$;
it does not supply the comparisons below.

\subsection{Exact edge laws and their entropy cost}
Let $\Omega$ be a finite nonempty alphabet, let $\kappa$ be a positive
probability law on it, and let $\lambda$ be a symmetric pair law with both
marginals $\kappa$.  Write $R_G=\kappa^{V(G)}$ and
\[
 D(\nu\Vert R)=\sum_x\nu(x)\log\frac{\nu(x)}{R(x)},
 \qquad I(\lambda)=D(\lambda\Vert\kappa^{\otimes2}),
\]
with natural logarithms and $0\log0=0$.  An exact edge witness on $G$
has singleton law $\kappa$ at every actual vertex and pair law $\lambda$
on every oriented actual edge.  Prescribing an averaged edge law is a
different problem.  The relation between entropy maximization and graph
density inequalities is developed in \citet{szegedy2015}; a bound for one
scalar density is not substituted here for exact marginal constraints.

\begin{theorem}[Potts edge witnesses]\label{theorem-potts-edge-entropy}
Let $G$ be a finite simple graph with $m$ edges, each of which has an
endpoint of degree at most three.  Let $q\ge2$, let $\kappa$ be uniform
on $\{1,\ldots,q\}$, and put
\[
 \lambda_s(x,y)=
 \begin{cases}s/q,&x=y,\\(1-s)/(q(q-1)),&x\ne y,\end{cases}
 \qquad 0\le s\le1.
\]
If $s<1/q$, assume additionally that $G$ is bipartite.  There is an
exact edge witness $\tau$ with
\[
 D(\tau\Vert R_G)\le m I_q(s),\qquad
 I_q(s)=s\log(qs)+(1-s)\log\frac{q(1-s)}{q-1}.
\]
Disconnected graphs and isolated vertices are included.
\end{theorem}

The theorem includes every bipartite $K_4$ subdivision and every uniform
Potts pair law.  Its antiferromagnetic part retains bipartiteness; it does
not extend to arbitrary pair tables by permutation symmetry.  The proof
uses a constrained finite Gibbs fit.  Two-colour conditioning supplies the
correlation sign, following the mechanism in \citet{csikvarilin2017}; a
conditional comparison at a low-degree endpoint makes every zero-coupling
edge fit.  A reverse-KL interpolation gives the sharp cost.  These steps,
including zero-support endpoints, are detailed in
Appendix~\ref{appendix-entropy-density}.

\begin{corollary}[Strict slack and full-table perturbations]
\label{corollary-potts-table-openness}
Under the preceding graph hypotheses, suppose $G$ has a cycle and
$0<s<1$, $s\ne1/q$.  The fitted Potts witness has
$D(\tau\Vert R_G)<mI_q(s)$.  There is a neighbourhood of
$(\kappa,\lambda_s)$, relative to compatible positive singleton and
symmetric pair laws, in which every input admits an exact witness with
$D(\tau\Vert\kappa^{V(G)})<mI(\lambda)$.
For $q\ge3$ the neighbourhood may be chosen inside the class of pairs
having a positive exchangeable triple extension.
\end{corollary}

This is a full-table neighbourhood: it includes nonuniform singleton laws
and non-Potts tables.  Its radius is not asserted to be uniform, and no
openness is inferred at $s=0,1$ or at independence.  At independence a
separate argument is available.

\begin{theorem}[Exact tables near independence]
\label{theorem-local-table-entropy}
Fix a finite simple bipartite graph $G$ containing a cycle, with girth
$g$, and a positive law $\kappa$ on an alphabet of size $q\ge2$.
For a symmetric pair law $\lambda$ with marginal $\kappa$, set
\[
 f(x,y)=\frac{\lambda(x,y)}{\kappa(x)\kappa(y)}-1,
 \qquad C_{xy}=\sqrt{\kappa(x)\kappa(y)}\,f(x,y),
 \qquad \rho=\|C\|_{\mathrm{op}}.
\]
There is a positive, explicitly specified radius $\delta(G,\kappa)$
such that $0<\rho\le\delta$ implies the existence of a full-support
exact witness with
\[
 D(\tau\Vert\kappa^{V(G)})\le |E(G)|I(\lambda)-\rho^g/8.
\]
At $\rho=0$ the independent law works.  For a forest, an exact witness
satisfies $D(\tau\Vert\kappa^{V(G)})=|E(G)|I(\lambda)$ by a tree
construction, without a smallness condition.
\end{theorem}

The local signed-density framework is already established
\citep{lovasz2010}.  Here a nearly fitted pairwise Gibbs seed is corrected
in all unary and edge-table coordinates.  The positive shortest even-cycle
term pays for that correction; the proof and its conservative constants
are in Appendix~\ref{appendix-entropy-density}.  Small scalar density
error alone would not supply the prescribed marginals.

\begin{proposition}[Cycle surplus for disjoint alphabet blocks]
\label{proposition-alphabet-block-surplus}
Let $G$ have $n$ vertices, $m$ edges and $c$ components, and write
$\beta=m-n+c$.  Partition the alphabet into nonempty blocks with masses
$\pi_i>0$.  In block $i$, let $\kappa_i$ and $\lambda_i$ be compatible
singleton and symmetric pair laws.  Suppose each connected component $C$
has an exact within-block witness of cost at most
$|E(C)|D(\lambda_i\Vert\kappa_i^{\otimes2})$.
Set $\kappa=\sum_i\pi_i\kappa_i$ and
$\lambda=\sum_i\pi_i\lambda_i$, with the pair law supported on
same-block pairs.  There is an exact witness $\tau_0$ with
\[
 D(\tau_0\Vert\kappa^{V(G)})\le mI(\lambda)-\beta H(\pi).
\]
If a known exact witness has slack $\Delta$, mixing its pair and joint
laws with $\varepsilon$ times their independent reference laws gives
a new exact witness with slack at least
$(1-\varepsilon)\Delta-mh(\varepsilon)$, where $h$ is binary entropy.
\end{proposition}

\begin{proof}
Choose one block independently in each connected component, then use its
within-block witness.  Disjoint supports give the exact component cost
$(|V(C)|-1)H(\pi)+\sum_i\pi_iD(\tau_{i,C}\Vert\kappa_i^{V(C)})$.
The pair cost is $I(\lambda)=H(\pi)+\sum_i\pi_iI(\lambda_i)$.
Summing proves the first bound, including isolated vertices.
For the mixture statement, joint KL convexity bounds the new joint cost
by $(1-\varepsilon)$ times the old cost.  The pair-level chain rule gives
\[
 (1-\varepsilon)I(\lambda_0)
   =I((1-\varepsilon)\lambda_0+\varepsilon\kappa^{\otimes2})
     +I(B;X,Y),
\]
where $B$ is the Bernoulli mixture label and $I(B;X,Y)\le h(\varepsilon)$.
Subtracting gives the claimed slack.
\end{proof}

The term $\beta H(\pi)$ permits a global block choice to cost entropy
while retaining a sharp edge budget.  It cannot be recovered by imposing
a separate sharp bound on every projected pair law.  For example, on the
six-vertex bipartite $K_4$ subdivision with path lengths
$(2,1,1,1,1,2)$ in the order $ab,ac,ad,bc,bd,cd$, take uniform
$\kappa$ on three symbols and
\begin{equation}\label{eq-three-colour-budget-transfer}
 \lambda=\frac1{5400}
 \begin{pmatrix}1787&7&6\\7&599&1194\\6&1194&600\end{pmatrix}.
\end{equation}
This positive table has a positive exchangeable triple extension.  An
explicit exact witness satisfies $8I(\lambda)-D(\tau\Vert\kappa^6)
>229/540$.  Its projected $\{1,2\}$-indicator pair is independent,
but its global indicators need not be; the full budget pays for them.
The construction and all positivity checks are given in
Appendix~\ref{appendix-entropy-density}.  This covers the displayed
table on the displayed graph, not arbitrary path lengths or all tables.
The graph $K_{3,3}-e$ already lies in the scalar-density class of
\citet{conlonfoxsudakov2012}.  Combined with the universal entropy
formulation in \citet[Corollary~9.1]{szegedy2015}, this already supplies
the sharp exact-pair entropy bound on this graph for every pair law.
The example records an explicit budget-transfer construction, not a new
graph or table family.

\begin{problem}[The unrestricted finite-alphabet question]
\label{problem-unrestricted-alphabet}
Suppose $\lambda$ is the pair marginal of an exchangeable law on
$\Omega^3$.  Does every actual simple bipartite $K_4$ subdivision $J$
admit an exact witness satisfying
$D(\tau\Vert\kappa^{V(J)})\le |E(J)|I(\lambda)$?
Zero entries in the pair and triple laws are allowed.
\end{problem}

The results above settle specified families, not this problem.  Exact
feasibility itself is easy: broadcast one pair sampled from $\lambda$
to the two bipartition classes.  This may cost too much entropy.  Even
for binary data, gluing stationary Markov bridges can impose incompatible
branch correlations.  Thus the remaining issue is a global sharp budget,
not merely an inverse fit or translation into Lean.

\subsection{Powers of one weighted regular host}
Let $X$ be finite, let $\mu_x>0$ with $\sum_x\mu_x=1$, and let
$W:X^2\to[0,1]$ be symmetric with
$\sum_y\mu_yW(x,y)=p>0$.  Composition and inner products use this
original measure.  Set $T=W/p$, $A=T^2$, $K_j=A^j$,
$\Pi(x,y)=1$ and $I(x,y)=\mathbf1_{x=y}/\mu_y$.
For six positive integers $a_e$, indexed by the edges of $K_4$, put
\[
 F(a)=\int_{X^4}\prod_{e=vw}K_{a_e}(x_v,x_w)\,d\mu^4.
\]
If $J(2a)$ replaces edge $e$ by a path of length $2a_e$, then
\[
 F(a)=\frac{t(J(2a),W)}{p^{2\sum_ea_e}},
\]
where $t$ denotes weighted homomorphism density.
This reduction concerns all-even subdivisions;
it does not cover every bipartite parity pattern.

\begin{theorem}[Two centered spectral levels]
\label{theorem-two-level-density}
Suppose, on the complete space orthogonal to constants,
$A$ has at most two eigenvalues.  Equivalently,
\[
 A=yI+(1-y)\Pi+(x-y)P,
 \qquad 0\le y\le x\le1,
\]
where $P$ is an orthogonal projection with $P1=0$.  Then
$F(a)\ge1$ for every choice of the six positive integers $a_e$.
Both spectral multiplicities are unrestricted.
\end{theorem}

The whole centered spectrum is meant: two nonzero eigenvalues together
with an additional zero complement give three levels and are not covered.
The proof places every power in the cone generated by $I$, $\Pi$ and
one common nonnegative positive-semidefinite Markov kernel.  Weighted
identity contractions reduce the cone expansion to a finite list of
quotients of $K_4$; their comparisons are proved in the appendix.
Single-kernel positive-spectrum inequalities are studied by
\citet{sidorenko2021}.  Neither that literature nor the cone argument
permits replacing one common host by six unrelated PSD kernels.

The appendix also shows strictness unless $A=\Pi$ when $|X|\ge2$.
For fixed powers this gives a neighbourhood of actual roots that may have
more than two spectral levels.  More precisely, for symmetric nonnegative
Markov roots $T,T_0$ on the same $(X,\mu)$, write
$N=\sum_ea_e$, $M_0=\|T_0\|_\infty$ and
$\varepsilon=\|T-T_0\|_\infty$, using maximum absolute kernel entries.
Markov averaging and telescoping give
\[
 |F(a;T)-F(a;T_0)|\le2N\varepsilon(M_0+\varepsilon)^5.
\]
Indeed, $\|T^j-T_0^j\|_\infty\le j\varepsilon$ and every positive
power is bounded by $M_0+\varepsilon$; telescope the six factors.
This elementary stability observation is host- and tuple-dependent,
not uniform compensation for arbitrary interacting spectra.

\begin{problem}[The unrestricted common-power question]
\label{problem-unrestricted-common-powers}
For the same actual weighted regular host, is
\[
 \int_{X^4}K_k(ab)K_{r+h}(ac)K_u(ad)
 K_r(bc)K_u(bd)K_\ell(cd)\,d\mu^4\ge1
\]
whenever $k\ge u\ge1$, $r\ge\ell\ge1$ and $h\ge1$,
without further spectral or mixing hypotheses?
\end{problem}

This question remains open here.  Individual centered terms and individual
spectral contributions can be negative even for actual Markov-square hosts; the full
density retains cycle, diamond and constant-mode surplus.  An arbitrary
PSD operator without a nonnegative Markov square root is not an admissible
counterexample.  These restrictions also explain why the two open problems
cannot be closed by a scalar calculation for a single test host.
The subdivision results of \citet[Theorems~1.3, 1.4 and 1.6]{imliliu2026}
provide important known cases: uniform even-theta substitution, nonuniform
substitution under divisibility conditions, and a clique subdivision with
one exceptional vertex.  We do not treat these cases as new results or
infer arbitrary nonuniform path lengths from them.

\section*{Scope of formal verification}

The classification and specified finite structural extensions have verified
Lean~4 implementations \citep{demoura2021,mathlib2020}.  The accepted results
include the canonical atoms, finite parameter and connected atom-count spectra,
classical-block identification, local decomposition products, actual-piece
accounting and supported-cover increments.  They use the original carriers
and retain their exact finiteness and structural hypotheses.

Verification of the extended manuscript is not complete.  In particular,
none of the entropy--density statements in
Section~\ref{sec:quantitative-cores} has an accepted Lean implementation.
Moreover,
fixed-atom profile realization, the full lattice spectra, maximal-chain and
height endpoints, the complete phase/equality interfaces and the new
coefficient bounds remain distinct obligations.  The arithmetic deficit
theorem supports the cyclic-budget calculation, but does not by itself
identify every remaining geometric or lattice endpoint.  The proofs in the
main text establish the manuscript conclusions; successful typechecking of
other statements cannot replace them.  Implementation details, exact exported
types, validation chronology and the remaining crosswalk are separated into
Appendix~\ref{appendix-formalization}.

\appendix
\section{Formalization interfaces and reproducibility}\label{appendix-formalization}

The implementation uses Lean~4 \citep{demoura2021} and the mathematical
library Mathlib \citep{mathlib2020}.  These references describe the proof
assistant and library, not the verification status of this manuscript.
Formal verification checks a precisely stated theorem; its correspondence
with a prose formulation still requires mathematical review
\citep{avigadharrison2014}.  Verified typechecking addresses a distinct trust
question \citep{carneiro2025}; no use of that external checker is claimed here.

The self-contained Lean~4 source is available in the public
\href{https://github.com/SamPetkov/Erdos593/blob/main/formalization/Erdos593SelfContained.lean}{\texttt{Erdos593} repository}.
The repository contains a classification implementation in
\href{https://github.com/SamPetkov/Erdos593/commit/6fd00a76064401f3f10aabef474f59d3c6ecd6bf}{commit \texttt{6fd00a7}},
whose Git timestamp is 15 July 2026.  A dated public validation record is
\href{https://github.com/SamPetkov/Erdos593/actions/runs/29682191161}{the successful GitHub Actions run of 19 July 2026},
at commit \texttt{88ce701}.  A Git timestamp alone does not establish the date
of public disclosure or priority.  Li's arXiv v2, submitted on 23 July 2026,
reports a separate formal verification of the broader results in his paper.
These records do not date either author's private work, and no shared code
or derivation between the two Lean developments is asserted.
For every finite triple system $F$, the exported theorems establish
\[
\begin{aligned}
\mathtt{F.IsObligatory}
&\Longleftrightarrow \mathtt{F.isolatedReduction.Intrinsic},\\
\mathtt{F.IsObligatory}
&\Longleftrightarrow \mathtt{Constructible\ F.isolatedReduction}.
\end{aligned}
\]
These exported theorems verify both directions of the classification,
including the isolated-vertex reduction.  Here ``finite'' describes the
system $F$ being classified.  The host triple systems are not assumed finite;
they are quantified within the documented ambient-universe convention of the
formalisation.

Further accepted Lean results cover the canonical-atom normal form and
minimal-generator classification, the finite parameter spectra, the connected atom-count
spectrum and its maximisers, the unrestricted classical-block equivalence,
the forward atomic cycle and even-theta boundary forms, and the supported
decomposition order isomorphism with local partition products.  A separately
validated extension has passed canonical and standalone replay and final
source reconciliation for the actual-piece/local-partition block-count
identity; these results retain their stated finiteness, structural and
feasible-parameter hypotheses.  The block-count checkpoint is now included
in the synchronized public and private repositories.  The abstract
capacity-two incidence-forest construction
has also been validated, but it does not yet formalise assembly of an arbitrary
fixed list of actual atoms.

The supported-subtype cover criterion, the unique local coordinate changed by
a cover, and the additive piece-deficit identity have now passed pinned
canonical and standalone Lean replay and final live source reconciliation.
This extension also verifies that coarsening along a cover reduces the actual
piece count by one and increases the deficit by one. These results use actual
original-edge pieces and local quotient carriers. The extension is included
in the public repository through
\href{https://github.com/SamPetkov/Erdos593/pull/51}{the reviewed cover/grading
integration}. Maximal-chain lengths and height endpoints
remain separate obligations.

Verification of the extended manuscript is therefore partial.  Remaining
interfaces include the maximal-chain grading and height endpoints, fixed-atom
profile realization and
canonical-label fidelity after assembly, the componentwise atom-count and full factorization-lattice
spectra, the all-maximal-block running-order statement, and sorted theta
parameter normalization.  The full indecomposability phase diagram,
decomposable spectrum and decomposable equality classification also remain
separate obligations.  The characteristic-polynomial recovery of the
attachment profile is not claimed as an additional Lean endpoint.  Exact
statement-by-statement coverage is tracked in the manuscript crosswalk.

The accepted atom-deficit theorem supplies the arithmetic identity and
positive-atom bound behind
Proposition~\ref{proposition-cyclic-cost-budget}, under its exact structural,
positive-rank and additive-count hypotheses.  Its identification with lattice
rank still uses the manuscript's grading argument.  The generating polynomial,
coefficient bounds and equality cases in
Corollary~\ref{corollary-sharp-piece-counts} are proved here but are not yet
additional Lean endpoints.  Counting by actual pieces uses the accepted
product and block-count interfaces, not an assumption about covers in an
ambient partition lattice; the fixed-atom sharpness assertion retains the
separate assembly-realization obligation.

The Lean development does not reproduce every intermediate prose lemma
verbatim.  In particular, its negative direction uses a shift-graph host in
place of the paper's high-odd-girth input, and its closure construction uses
a singular-safe layering rather than the displayed cofinality-indexed chain.
These are distinct proof routes to the classification.  Kernel validation
supplements rather than replaces mathematical review; a full external audit
of the extended development remains pending.

% Technical proof fragment; main statements carry the theorem labels.
% Source: universal-round1-20261009, independently reviewed analytical proofs.
% This fragment makes no claim of Lean validation or unrestricted coverage.
\section{Exact marginal entropy and common-power densities}
\label{appendix-entropy-density}

All logarithms in this appendix are natural.  If $R$ is a strictly positive
probability law on a finite set, we use
$D(\nu\Vert R)=\sum_x\nu(x)\log(\nu(x)/R(x))$, with $0\log0=0$.
Every marginal below is attached to an actual vertex or edge; it is not an
average over edges.  The arguments establish the restricted statements in
Theorems~\ref{theorem-potts-edge-entropy}, \ref{theorem-local-table-entropy}
and~\ref{theorem-two-level-density}.  They do not resolve the arbitrary-table
entropy question or the unrestricted common-power inequality.

\subsection{Uniform Potts marginals}

\begin{proof}[Proof of Theorem~\ref{theorem-potts-edge-entropy}]
Write $G=(V,E)$, $m=|E|$, and let $R$ be the uniform product law on
$\{1,\ldots,q\}^{V}$.  The one-symbol case is immediate, so assume $q\ge2$.
For $0<s<1$, put
\[
 s_0=\frac1q,\qquad
 \theta(s)=\log\frac{(q-1)s}{1-s},\qquad
 I_q(s)=s\log(qs)+(1-s)\log\frac{q(1-s)}{q-1}.
\]
For an interaction vector $J=(J_e)_{e\in E}$, define
\begin{align*}
 Z(J)=\mathbb E_R\exp\!\left(\sum_{uv\in E}J_{uv}
                    \mathbf1_{\{X_u=X_v\}}\right),&\\
 \frac{d\nu_J}{dR}=Z(J)^{-1}
       \exp\!\left(\sum_{uv\in E}J_{uv}\mathbf1_{\{X_u=X_v\}}\right).&
\end{align*}
Minimize $\log Z(J)-s\sum_eJ_e$ on $[0,\theta(s)]^E$ if $s>s_0$,
and on $[\theta(s),0]^E$ if $s<s_0$.  A minimizer exists and has full
support.  Simultaneous permutations of the $q$ colours preserve its law.
Thus each singleton is uniform and every ordered pair is determined by its
agreement probability: all diagonal entries are equal, as are all
off-diagonal entries.  It remains to fit every edge agreement to $s$.

\smallskip\noindent\emph{The removed-edge sign.}
For nonnegative interactions, any two distinct vertices have agreement
probability at least $1/q$.  For nonpositive interactions on a bipartite
graph, vertices in opposite parts have agreement probability at most $1/q$.
Here is a finite proof of both assertions.  Condition on the set of vertices
whose colours lie in $\{1,2\}$ and on every colour outside this set.  An edge
from the set to its complement always has unequal colours and introduces no
field.  The remaining law is a zero-field Ising model with couplings $J_e/2$.
For nonnegative couplings, the high-temperature expansion expresses its
correlation numerator as a sum over edge sets with boundary $\{u,v\}$,
with nonnegative products of $\tanh(J_e/2)$; its denominator is a positive
sum over even edge sets.  The correlation is therefore nonnegative.
For nonpositive interactions on a bipartite graph, reversing the spins on
one part gives the preceding nonnegative-coupling model, and reverses the
correlation between opposite parts.  Averaging the conditioning compares
the probabilities of colours $(1,1)$ and $(1,2)$.  By colour symmetry these
are respectively $a/q$ and $(1-a)/(q(q-1))$, where $a$ is the agreement
probability.  Their comparison gives the asserted sign.  This is the
weighted two-colour conditioning mechanism used in the proper-colouring
setting by \citet{csikvarilin2017}.

\smallskip\noindent\emph{Nonzero interactions fit.}
The derivative in coordinate $e=uv$ is $a_e-s$, where
$a_e=\nu_J(X_u=X_v)$.  An interior coordinate therefore fits.
At the far face $J_e=\theta(s)$, remove this interaction and call the
resulting agreement probability $a_0$.  Restoring it multiplies agreement
odds by $\exp\theta(s)$, so the new agreement is
\[
 \frac{e^{\theta(s)}a_0}{1+(e^{\theta(s)}-1)a_0}.
\]
The removed-edge sign shows that this is at least $s$ in the positive
regime and at most $s$ in the negative regime.  At the upper face the
derivative is nonpositive; at the lower face it is nonnegative.  These
Karush--Kuhn--Tucker inequalities give the reverse comparison in each
regime.  Hence every nonzero interaction fits, including a far-face one.

\smallskip\noindent\emph{Zero interactions fit.}
Suppose $J_{vw}=0$, and choose the endpoint $v$ with degree at most three.
There are at most two other neighbours with nonzero interactions.  If these
are $a,b$, write their interactions as $J_1,J_2$.  Conditioning on all
colours except $X_v$ gives
\[
 \nu_J(X_v=X_a)-\nu_J(X_v=X_b)
 =\nu_J(X_a\ne X_b)
       \frac{e^{J_1}-e^{J_2}}{q-2+e^{J_1}+e^{J_2}}.
\]
The left side is zero because both nonzero incident interactions fit.
Full support makes $\nu_J(X_a\ne X_b)>0$, so $J_1=J_2$.
For a positive common interaction, the colour of either active neighbour
maximizes the conditional probability of $X_v$, whether the two neighbour
colours agree or differ.  Consequently the agreement on $vw$ is at most
that on $va$, namely $s$.  The lower-face derivative inequality gives
agreement at least $s$.  For a negative common interaction, neighbour
colours minimize the conditional probability and both comparisons reverse.
With one active neighbour the same maximum or minimum comparison applies.
With none, agreement on $vw$ is $1/q$, contrary to the relevant zero-face
inequality when $s\ne s_0$.  Thus every original edge fits exactly.  Colour
symmetry now gives the whole prescribed Potts table, not just agreement.
Isolated vertices remain independent uniform variables.

\smallskip\noindent\emph{The entropy bound.}
For each interior $u$ on the segment from $s_0$ to $s$, choose a fitted law
$\nu_u$ as above, and take $\nu_{s_0}=R$.  Set
$C_u=D(\nu_u\Vert R)$ and $S_u=\sum_eJ_e(u)$.  Direct substitution gives
\begin{equation}\label{eq-potts-reverse-kl}
 D(\nu_v\Vert\nu_u)=C_v-C_u+(u-v)S_u\ge0.
\end{equation}
If $s_0\le v<u\le s$, then
$C_u-C_v\le m(u-v)\theta(u)$.
If $s\le u<v\le s_0$, then $S_u\ge m\theta(u)$ and the negative factor
$u-v$ again gives $C_u-C_v\le m(u-v)\theta(u)$.
Fix the target law $\nu_s$, telescope these inequalities along a finite
mesh from $s_0$ to $s$, and let its mesh tend to zero.  The oriented
Riemann sums yield
\[
 C_s\le m\int_{s_0}^{s}\theta(u)\,du=mI_q(s).
\]
This step requires neither continuity nor uniqueness of the chosen fitted
family.  In particular, the order of the two laws in
\eqref{eq-potts-reverse-kl} is essential in the negative regime.

At $s=s_0$, use $R$.  At $s=0$ or $s=1$, apply the interior result to
$s_\varepsilon=(1-\varepsilon)s+\varepsilon/q$ and take a convergent
subsequence in the finite joint simplex.  Marginalization and divergence
relative to the positive reference $R$ are continuous, including at zero
joint masses.  The limit gives all endpoint marginals and the bound.
The assertion is vacuous on edges when $E$ is empty, and the product law
then has zero cost.
\end{proof}

\subsection{Cyclic slack and perturbation of the full pair table}

\begin{proof}[Proof of Corollary~\ref{corollary-potts-table-openness}]
Keep the notation of the preceding proof, with $0<s<1$, $s\ne s_0$, and
with at least one cycle.  First, every minimizing interaction is nonzero.
Indeed, at a hypothetical zero edge the preceding neighbour-colour
comparison is strict on the positive-probability event
$X_a=X_b\ne X_w$ when there are two active neighbours, or $X_a\ne X_w$
when there is one.  Simplicity ensures that these are distinct actual
vertices.  Averaging contradicts the relevant zero-face inequality.
The case of no active neighbour was already excluded.

If $e=uv$ lies on a cycle, removing its interaction leaves an alternative
$u$--$v$ path with nonzero interactions.  In the two-colour conditioning,
every conditional correlation has the weak sign established above.
The event that all vertices have colours in $\{1,2\}$ has positive
probability.  On this event the Ising high-temperature numerator contains
the strictly positive contribution of the alternative path, after a
bipartition spin reversal in the negative regime.  The removed-edge
agreement is therefore strictly above $1/q$ in the positive regime and
strictly below $1/q$ in the negative regime.  Reinsertion at the interaction
cap would strictly overshoot $s$, contradicting the cap inequality.  Hence
\[
 0<J_e(s)<\theta(s)\quad(s>s_0),\qquad
 \theta(s)<J_e(s)<0\quad(s<s_0)
\]
on every cycle edge.  Bridges need not be strictly inside their cap.

The covariance matrix of the statistics
$h_e=\mathbf1_{\{X_u=X_v\}}$ is positive definite under every full-support
law.  To see this, restrict a putatively constant linear combination to
two colours and write $h_e=(1+\sigma_u\sigma_v)/2$.  The distinct
simple-edge Walsh monomials are linearly independent modulo constants,
so every coefficient must vanish.  Thus $\log Z$ has positive-definite
Hessian and an injective gradient.  The analytic implicit function theorem
supplies a unique locally analytic fitted vector $J(s)$; uniqueness patches
these local inverses along the applicable interval.  That interval is
$0<s<1$ for a bipartite graph, and $s_0\le s<1$ in the general positive
regime.  No negative-regime fit is asserted for a nonbipartite graph.
The zero vector is the unique fit at $s_0$.

For the exact fit,
\[
 C(s)=s\sum_eJ_e(s)-\log Z(J(s)),\qquad C'(s)=\sum_eJ_e(s).
\]
The terms involving $J'(s)$ cancel because every edge agreement is $s$.
The cap bounds imply $J(s)\to0$ at $s_0$, and hence
\begin{equation}\label{eq-potts-strict-gap}
 \Delta_s:=mI_q(s)-C(s)
 =\int_{s_0}^{s}\left(m\theta(u)-\sum_eJ_e(u)\right)du>0.
\end{equation}
Above $s_0$ the integrand is positive because a cycle edge is strictly
inside its cap.  Below $s_0$ the integrand is negative and the orientation
of the integral reverses.  This proves strictness in both regimes.

To perturb the whole table, fix a mean-zero orthonormal basis
$f_1,\ldots,f_{q-1}$ for the uniform singleton law $\kappa_0$.
Use every unary statistic $f_i(X_v)$ and, for a fixed orientation of each
edge, every statistic $f_i(X_u)f_j(X_v)$.
These statistics are orthonormal under $\kappa_0^V$, since distinct
supports leave an unmatched centred factor.  They are therefore affinely
independent modulo constants on the joint configuration space, and their
covariance under the full-support law $\nu_s$ is positive definite.
The expectation map for exponential tilts of $\nu_s$ is locally invertible.
It fits every sufficiently nearby compatible singleton and symmetric pair
table: the unary and pair-product basis reconstructs every full actual-edge
table.  The basis is fixed while the target singleton law varies.
The tilted joint law depends continuously on these moments.  Both
$D(\tau\Vert\kappa^V)$ and $mD(\lambda\Vert\kappa^{\otimes2})$ are
continuous for positive nearby inputs.  By \eqref{eq-potts-strict-gap},
their difference remains negative after shrinking the neighbourhood.
This gives an open family of non-Potts tables and nonuniform singleton laws,
but not coverage of the entire compatible-table domain.

For $q\ge3$, this neighbourhood can also remain within the positive
exchangeable-triple cone.  Let $\eta_{\mathrm{iid}}$ be the uniform product
triple, $\eta_{\mathrm{dist}}$ the uniform ordered all-distinct triple,
and $\eta_{\mathrm{eq}}$ the uniform all-equal triple.  A positive extension
of the base table is
\[
 \eta_s=qs\eta_{\mathrm{iid}}+(1-qs)\eta_{\mathrm{dist}}
 \quad(0<s<s_0),
\]
and
\[
 \eta_s=\frac{1-s}{1-s_0}\eta_{\mathrm{iid}}
       +\left(1-\frac{1-s}{1-s_0}\right)\eta_{\mathrm{eq}}
 \quad(s_0<s<1).
\]
For a nearby symmetric probability-table perturbation
$\lambda'=\lambda_s+\delta$, set $d_x=\sum_y\delta(x,y)$ and define
\begin{align*}
 H(x,y,z)={}&\delta(x,y)\kappa_0(z)+\delta(x,z)\kappa_0(y)
             +\delta(y,z)\kappa_0(x)\\
 &-d_x\kappa_0(y)\kappa_0(z)-d_y\kappa_0(x)\kappa_0(z)
             -d_z\kappa_0(x)\kappa_0(y).
\end{align*}
This is symmetric, has total mass zero, and its $(x,y)$ marginal is
$\delta(x,y)$: the two row-drift terms cancel and $\sum_zd_z=0$.
Thus $\eta_s+H$ extends $\lambda'$ and has singleton
$\kappa_0+d$.  It stays strictly positive for sufficiently small
perturbations.  This last assertion is deliberately restricted to $q\ge3$;
the binary triple-extension cone has additional boundary restrictions.
\end{proof}

\subsection{An explicit neighbourhood of independence}

\begin{proof}[Proof of Theorem~\ref{theorem-local-table-entropy}]
Let $G$ be finite, simple and bipartite, let $n=|V|$, $m=|E|$, and let
$\kappa$ be a positive law on a $q$-symbol alphabet, $q\ge2$.
For a symmetric table $\lambda$ with marginals $\kappa$, put
\begin{gather*}
 a=\min_x\kappa(x),\qquad
 f(x,y)=\frac{\lambda(x,y)}{\kappa(x)\kappa(y)}-1,\\
 C_{xy}=\sqrt{\kappa(x)}f(x,y)\sqrt{\kappa(y)},\qquad
 \rho=\|C\|_{\mathrm{op}}.
\end{gather*}
If $G$ is a forest, root each component, draw its root from $\kappa$, and
draw each child using transition $\lambda(x,y)/\kappa(x)$, independently
conditional on its parent.  This has the required marginals and divergence
$mD(\lambda\Vert\kappa^{\otimes2})$ by the chain rule, including zero
pair entries.  Isolated vertices contribute zero.

Suppose now that $G$ has even girth $g\ge4$.  The following constants specify
the radius used in the theorem:
\begin{gather*}
 d=n(q-1)+m(q-1)^2,\quad A=2^ma^{-m},\quad M=a^{-1},
\\ B=\sqrt d\,(M+1)A,\quad L=\sqrt d\,M,\quad r=\frac{\log2}{4L},\\
 \begin{aligned}
 \delta_0=\min\Bigl\{&1,\frac a3,\frac1{2A},\frac{a\log2}{8m},\\
       &\left(\frac r{8B}\right)^{1/(g-1)},
       \left(\frac1{8B^2}\right)^{1/(g-2)}\Bigr\}.
 \end{aligned}
\end{gather*}
Assume $0<\rho\le\delta_0$.  Every matrix entry of $C$ has magnitude at
most $\rho$, so $|f(x,y)|\le\rho/a\le1/3$.
With $R=\kappa^V$, take the full-support seed
\[
 w(X)=\prod_{uv\in E}(1+f(X_u,X_v)),\qquad
 Z=\mathbb E_Rw,\qquad \nu_0=Rw/Z.
\]
This seed is not assumed to have the target marginals.

Expand $Z$ over actual edge subsets.  A subset with a leaf contributes
zero, because each $\kappa$-weighted row sum of $f$ is zero.  A surviving
nonempty subset has at least $g$ edges, and a surviving $g$-edge subset is
exactly a shortest cycle.  Its contribution is $\operatorname{Tr}(C^g)$,
which is at least $\rho^g$ because $g$ is even and $C$ is real symmetric.
The absolute total of all terms with at least $g+1$ edges is at most
$A\rho^{g+1}$: a $k$-edge term is bounded by $(\rho/a)^k$, and there are
at most $2^m$ terms.  Similarly $|Z-1|\le A\rho^g$.  Therefore
\begin{equation}\label{eq-local-partition-gain}
 Z\ge1+\rho^g-A\rho^{g+1}\ge1+\frac{\rho^g}{2},\qquad
 \log Z\ge\frac{\rho^g}{4}.
\end{equation}
The last inequality uses $\rho\le1$.  The even-cycle mechanism is the
finite weighted version of the local signed-density argument in
\citet{lovasz2010}.  We still need to correct the seed to exact marginals;
a scalar partition bound alone does not do so.

Choose a $\kappa$-orthonormal mean-zero basis
$\xi_1,\ldots,\xi_{q-1}$.  Let $T=(T_1,\ldots,T_d)$ contain all unary
statistics $\xi_i(X_v)$ and all oriented actual-edge products
$\xi_i(X_u)\xi_j(X_v)$.  Under $R$ they are mean zero and orthonormal,
including on distinct edges sharing a vertex.  Each has magnitude at most
$M$, since $|\xi_i(x)|\le1/\sqrt{\kappa(x)}$.
The seed potential has the complete gauge decomposition
\[
 h_0=\sum_{uv\in E}\log(1+f(X_u,X_v))=\theta_0\cdot T+c.
\]
Let $b$ be the desired moment vector: its unary coordinates vanish and its
edge coordinates are $\sum_{xy}\lambda(x,y)\xi_i(x)\xi_j(y)$.
Each of the latter has magnitude at most $\rho$, by the definition of $C$.
Matching $b$ fixes every singleton and every full edge table, because
$\{1,\xi_1,\ldots,\xi_{q-1}\}$ is a basis of all functions on the alphabet.

We bound the seed residual explicitly.  In a unary-inserted edge-subset
expansion, every unmarked leaf kills the term.  A nonempty forest has at
least two leaves, so every surviving term contains a cycle and at least
$g$ edges.  For an edge statistic marked at adjacent $u,v$, a surviving
forest can only be a single $u$--$v$ path.  The one-edge path contributes
exactly the desired coordinate of $b$; every other path has at least
$g-1$ edges, as it closes with $uv$ to a cycle.  Cyclic terms have at least
$g$ edges.  Thus each unnormalized edge moment is $b_i+E_i$ with
$|E_i|\le MA\rho^{g-1}$.  Division by $Z\ge1$ adds at most
$|b_i||Z-1|\le A\rho^{g+1}$.  The unary bound is no larger.  Consequently
\begin{equation}\label{eq-local-marginal-residual}
 \|\mathbb E_{\nu_0}T-b\|_2\le B\rho^{g-1}.
\end{equation}

Consider the convex function
\[
 F(\theta)=\log\mathbb E_R e^{\theta\cdot T}-\theta\cdot b.
\]
Its gradient at $\theta_0$ is the residual in
\eqref{eq-local-marginal-residual}, and its Hessian is the covariance of
$T$ under its Gibbs law.  The range of $h_0$ is at most
$4m\rho/a\le(\log2)/2$, using $|\log(1+f)|\le2|f|$.
A displacement of Euclidean norm at most $r$ adds potential range at most
$2Lr=(\log2)/2$.  Throughout the closed radius-$r$ ball about $\theta_0$,
the Gibbs density relative to $R$ is therefore at least $1/2$.
For every $z\in\mathbb R^d$,
\[
 \operatorname{Var}_{\nu_\theta}(z\cdot T)
 =\inf_c\mathbb E_{\nu_\theta}(z\cdot T-c)^2
 \ge\frac12\inf_c\mathbb E_R(z\cdot T-c)^2
 =\frac12\|z\|_2^2.
\]
Hence $\nabla^2F\ge\tfrac12I$ on that ball.
Write $G_0=\|\nabla F(\theta_0)\|_2\le B\rho^{g-1}\le r/8$.
Integration along a segment gives
\[
 F(\theta_0+z)\ge F(\theta_0)+\nabla F(\theta_0)\cdot z
                              +\frac14\|z\|_2^2.
\]
On the radius-$r$ boundary this exceeds $F(\theta_0)$ by at least $r^2/8$.
A minimum on the closed ball is therefore interior, has zero gradient,
and is a global minimum by convexity.  Its Gibbs law $\tau$ has full
support and fits $b$ exactly.  The same inequality gives
$F(\theta_0)-F(\theta_*)\le G_0^2$.

Put $I=D(\lambda\Vert\kappa^{\otimes2})$.
The target expectation of $h_0$ is $mI$, so the gauge constant cancels to
give $F(\theta_0)=\log Z-mI$.  Since the fitted law has moment vector $b$,
$D(\tau\Vert R)=-F(\theta_*)$.  Combining these facts yields
\[
 \begin{split}
 D(\tau\Vert R)
 &\le mI-\frac{\rho^g}{4}+B^2\rho^{2g-2}\\
 &\le mI-\frac{\rho^g}{8},
 \end{split}
\]
where the final step uses $\rho^{g-2}\le1/(8B^2)$.
At $\rho=0$, the table is $\kappa^{\otimes2}$ and the product joint works.
If a triple extension is required, it is supplied in this smallness regime
by the exchangeable probability law
\[
 \eta(x,y,z)=\kappa(x)\kappa(y)\kappa(z)
          [1+f(x,y)+f(x,z)+f(y,z)].
\]
Its bracket is nonnegative, its total mass is one, and its pair table is
$\lambda$.  No triple-extension hypothesis was used to construct $\tau$.
The radius depends on $G,q$ and the smallest singleton mass; the argument
does not cover arbitrary tables far from independence.
\end{proof}

\subsection*{The three-colour budget-transfer example}
Here we verify the assertion following
Proposition~\ref{proposition-alphabet-block-surplus} without a numerical
entropy estimate.  Let $G=K_{3,3}-e$, with the actual subdivision and table
in \eqref{eq-three-colour-budget-transfer}, and let $R$ be uniform on its
$3^6=729$ colourings.  Put
\[
 \lambda_0=\begin{pmatrix}1/3&0&0\\0&1/9&2/9\\0&2/9&1/9\end{pmatrix}.
\]
Its blocks have masses $1/3,2/3$.  The second block has uniform binary
singletons and agreement $1/3$.  The Potts proof supplies a Gibbs witness
$\tau_B$ with interactions in $[-\log2,0]$, hence every one of its
$2^6$ configurations has probability at least $2^{-14}$.
Choose all colours equal to $1$ with probability $1/3$, and choose
$\tau_B$ on $\{2,3\}$ with probability $2/3$; call the resulting law
$\tau_0$.  The block identity gives slack at least
$3H(1/3,2/3)$.  Set $\nu=(99/100)\tau_0+(1/100)R$ and
$\lambda_b=(99/100)\lambda_0+(1/100)\kappa^{\otimes2}$.
Then $8I(\lambda_b)-D(\nu\Vert R)\ge
(99/100)3H(1/3,2/3)-8h(1/100)$.

Let $D_{12}=D_{21}=1/5400$, $D_{11}=D_{22}=-1/5400$, with other
entries zero.  Its row sums vanish and $\lambda=\lambda_b+D$.
For $\psi=(1,-1,0)$, define
\[
 g(x,y)=\frac{D_{xy}}{\kappa_x\kappa_y}
       =-\frac{\psi(x)\psi(y)}{600},\qquad
 \tau(X)=\nu(X)+R(X)\sum_{uv\in E(G)}g(X_u,X_v).
\]
Every singleton correction vanishes.  At any actual edge, the direct
correction is $D$ and every other edge correction vanishes after integrating
a centred endpoint.  Thus total mass, singletons and every full edge table
are exact.  If a colouring contains colour $1$ and another colour, its
nontrivial vertex cut has at least two edges because $G$ is bridgeless.
Consequently $\sum_{uv}\psi(X_u)\psi(X_v)\le6$, so the negative
correction is paid by $R/100$.  With no colour $1$, the extra lower bound
$(99/100)(2/3)2^{-14}=33/819200>1/218700$ pays the possible deficit
$R(8/600-1/100)$.  The all-$1$ colouring has mass at least $33/100$
before correction and is also nonnegative afterwards.  These cases exhaust
all configurations, so $\tau$ is an actual probability law.

The pair correction has total variation $d_P=1/2700$, and joint variation
is at most $d_V=2/675$.  Maximal coupling gives
$|H(P)-H(Q)|\le h(d)+d\log(N-1)$ on $N$ states; the same bound
controls KL changes relative to a uniform law.  Its right side is increasing
for $0\le d\le1-1/N$.  Both $d_P$ and the upper bound $d_V$ lie in
their respective intervals, so substitution of the joint upper bound is valid.
Applying it to $N=9,729$
therefore gives slack at least
\[
 \frac{99}{100}3H(1/3,2/3)-8h(1/100)
 -8[h(d_P)+d_P\log8]-[h(d_V)+d_V\log728].
\]
Use $H(1/3,2/3)>\log(3/2)>1/3$ and
$h(t)\le t\log(1/t)+t$.  The elementary bound $e>8/3$ implies
$\log100<5$, $\log8<3$, $\log2700<9$, $\log728<8$ and
$\log(675/2)<7$.  The slack is strictly greater than
\[
 \frac{99}{100}-\frac{48}{100}-\frac{104}{2700}-\frac{32}{675}
 =\frac{229}{540}.
\]
For triple feasibility, let $\eta_0$ put mass $1/3$ on $111$ and mass
$1/3$ uniformly on each permutation orbit $223$ and $233$.
Its pair is $\lambda_0$.  The exchangeable law
\[
 \eta(x,y,z)=\frac{99}{100}\eta_0(x,y,z)
       +\frac1{100}\kappa_x\kappa_y\kappa_z
       +\frac{D_{xy}+D_{xz}+D_{yz}}3
\]
has pair $\lambda$ and is positive: its IID contribution is $1/2700$
and its correction is no smaller than $-1/5400$.
Finally, the projected $\{1,2\}$ indicator has pair probability $4/9$
for two ones and singleton probability $2/3$, hence an independent pair.
The construction makes no demand that its joint indicators be independent.

\subsection{A common-kernel cone on four vertices}

Let $(X,\mu)$ be a nonempty finite probability space with $\mu_x>0$.
Kernels act by $(Sf)(x)=\sum_y\mu_yS(x,y)f(y)$; all compositions, traces
and inner products below refer to this original measure.  Define
$I(x,y)=\mathbf1_{\{x=y\}}/\mu_y$ and $\Pi(x,y)=1$.
An operator power $S^j$ is not an entrywise power.

\begin{lemma}[Common-kernel cone]\label{lemma-common-kernel-cone}
Let $S$ be one symmetric, entrywise nonnegative, positive-semidefinite
Markov kernel on $(X,\mu)$.  For each of the six edges of $K_4$, choose
$L_e$ independently in $\operatorname{conv}\{I,\Pi,S\}$.
Then
\[
 \mathbb E_{\mu^4}\prod_{uv\in E(K_4)}L_{uv}(X_u,X_v)\ge1.
\]
The same actual kernel $S$ is used at all six edges.
\end{lemma}

\begin{proof}
Since $S1=1$ and $S$ is positive semidefinite, $S-\Pi$ is positive
semidefinite.  Thus $S(x,x)\ge1$.  Every eigenvalue of $S$ is
nonnegative and the constant eigenvalue is one, so
$\operatorname{Tr}(S^j)\ge1$ for every positive integer $j$.
The entrywise square $R(x,y)=S(x,y)^2$ is positive semidefinite as well.
Indeed, with $D=\operatorname{diag}(\mu)$ and $B=D^{1/2}SD^{1/2}$,
its weighted matrix is $D^{-1/2}(B\circ B)D^{-1/2}$; apply the Schur
product theorem.  We use that the trace of a product of two positive
semidefinite operators is nonnegative, but never assume a corresponding
sign for a product of three.

Expand the six convex combinations.  Their $3^6$ elementary assignments
of $I,S,\Pi$ have nonnegative coefficients summing to one.
Fix one such assignment.  In its $I$-edge graph choose a spanning forest
with the same components.  On the equality support of that forest, each
deleted cyclic $I$ edge is $1/\mu_x\ge1$; off its support the forest
product is zero.  All other factors are nonnegative, so deleting these
cyclic edges lowers the density.  Contract the remaining $I$ trees.
A tree with $t$ vertices contributes $\mu_x^t\mu_x^{-(t-1)}=\mu_x$.
Every quotient vertex therefore retains the original measure $\mu$.
Delete $\Pi$ factors and the resulting $S$ loops, which are at least one.
Parallel $S$ edges must be retained.  We now check all quotient types.

\smallskip\noindent\emph{Four quotient vertices.}
There was no contraction.  The retained simple graph is a subgraph of
$K_4$.  Ordinary leaves integrate to one using $S1=1$.
Its leaf-stripped core is empty, a triangle, a four-cycle, a diamond, or
$K_4$.  The first three have densities $1$, $\operatorname{Tr}(S^3)$,
and $\operatorname{Tr}(S^4)$ respectively.
Writing $S^2$ for operator composition, the diamond density is
\[
 \mathbb E_{c,d}S(c,d)[S^2(c,d)]^2
 \ge[\mathbb E_{c,d}S(c,d)S^2(c,d)]^2
 =[\operatorname{Tr}(S^3)]^2\ge1,
\]
by Jensen under the probability measure
$\mu_c\mu_dS(c,d)$.  For $K_4$, set $f_{cd}(x)=S(x,c)S(x,d)$.
Its density is
\[
 \mathbb E_{c,d}S(c,d)\langle f_{cd},Sf_{cd}\rangle
 \ge\mathbb E_{c,d}S(c,d)(\mathbb E_xf_{cd}(x))^2,
\]
because $S-\Pi$ is positive semidefinite and $S(c,d)\ge0$.
The last expression is the diamond density.
This finite argument also makes explicit the complete-graph
positive-kernel bound discussed in \citet{sidorenko2021}.

\smallskip\noindent\emph{Three quotient vertices.}
The original partition sizes are $(2,1,1)$; write its edge multiplicities
as $(s,t,v)\in\{0,1,2\}^2\times\{0,1\}$, where the two first edges
meet the merged vertex.  This gives all eighteen possibilities, including
the empty graph.  If $v=0$, the density is
\[
 \mathbb E_a\left(\mathbb E_bS(a,b)^s\right)
              \left(\mathbb E_cS(a,c)^t\right)\ge1,
\]
because each bracket is one for exponent zero or one and at least one
for exponent two by rowwise Jensen.  If $v=1$ and an arm is absent,
integrating an ordinary leaf leaves at most a double edge, bounded by
Jensen.  The remaining possibilities are $(1,1,1)$, $(2,1,1)$,
$(1,2,1)$ and $(2,2,1)$.  The first is a triangle.  Using the entrywise
square $R$, either middle case has density
\[
 \operatorname{Tr}(RS^2)\ge\operatorname{Tr}(R\Pi)
       =\mathbb E_{x,y}S(x,y)^2\ge1,
\]
since $R$ and $S^2-\Pi$ are positive semidefinite.  The last case has
density
\[
 \operatorname{Tr}(SR^2)\ge\operatorname{Tr}(\Pi R^2)
       =\|R1\|_{L^2(\mu)}^2
       \ge[\mathbb E_{x,y}S(x,y)^2]^2\ge1.
\]
Here $R^2$ is operator composition, not a fourth entrywise power.

\smallskip\noindent\emph{Two or one quotient vertices.}
Two vertices have partition sizes $(3,1)$ or $(2,2)$, hence at most four
parallel edges.  Their density is $\mathbb E S(x,y)^j\ge1$ by Jensen
for $j\ge1$; for $j=0$ it is one.  One quotient vertex has only the
already deleted loops and also has density one.
Every elementary assignment is therefore bounded below by one, proving
the convex-mixture assertion.  Zeros in $S$, unequal weights, reducibility,
and repeated eigenvalues cause no divisions or changes of measure.
\end{proof}

\begin{proof}[Proof of Theorem~\ref{theorem-two-level-density}]
Let $P$ be an orthogonal projection on $L^2(\mu)$ with $P1=0$, and assume
that the actual nonnegative kernel is
\[
 A=yI+(1-y)\Pi+(x-y)P,\qquad 0\le y\le x\le1.
\]
The mutually orthogonal projections $\Pi,P,I-\Pi-P$ give
\[
 A^j=y^jI+(1-y^j)\Pi+(x^j-y^j)P\qquad(j\ge1).
\]
If $x>y$ and $y<1$, put
$\alpha=(x-y)/(1-y)$ and $S=\Pi+\alpha P$.
This is symmetric, positive semidefinite and Markov.  For distinct
$u,v$, $S(u,v)=A(u,v)/(1-y)\ge0$, while
$S(u,u)=1+\alpha P(u,u)\ge1$ because $P$ is positive semidefinite.
Thus $S$ is an actual nonnegative kernel, not merely a formal spectral
projection.  Write
\[
 a_j=y^j,\qquad b_j=\frac{(1-y)(x^j-y^j)}{x-y},\qquad
 c_j=1-a_j-b_j.
\]
Then $A^j=a_jI+b_jS+c_j\Pi$, and the coefficients are nonnegative:
\[
 \frac{x^j-y^j}{x-y}
   =\sum_{i=0}^{j-1}x^{j-1-i}y^i
   \le\sum_{i=0}^{j-1}y^i=\frac{1-y^j}{1-y}.
\]
Apply Lemma~\ref{lemma-common-kernel-cone} to the six positive exponents,
independently.  No ordering of them and no rank-one restriction on $P$
is needed.

If $x=y<1$, use $A^j=y^jI+(1-y^j)\Pi$ directly.  If $y=1$, then
$x=1$ and $A=I$; its six-edge density is
$\sum_z\mu_z^{-2}\ge1$.  If $y=0$, the preceding argument includes
the zero complementary band, and when $x=0$ the kernel is simply $\Pi$.
These cases avoid every potentially zero denominator.
For $|X|\ge2$, equality holds precisely when $A=\Pi$.
When $y>0$, the all-$I$ assignment has positive weight and density
$\sum_z\mu_z^{-2}>1$.  When $y=0$ and $A\ne\Pi$, the all-$S$
assignment with $S=A$ has positive weight, and
\[
 t(K_4,A)\ge[\operatorname{Tr}(A^3)]^2
           =[1+\operatorname{rank}(P)x^3]^2>1.
\]
All remaining elementary terms are at least one.

The representation is a condition on the \emph{complete} centred spectrum:
when $x,y>0$, an additional centred zero eigenspace is not included.
The same-base-kernel condition of the lemma is also essential; the proof
does not assert a bound for unrelated, even commuting, nonnegative
positive-semidefinite kernels.  Nor does this restricted density result
give a universal exact-edge entropy witness; the universal entropy
formulation and its graph-density relation require their full quantifiers
\citep{szegedy2015}.
\end{proof}

\section*{Acknowledgments}

The author thanks Tom de Groot for his advice on revising the manuscript for
greater clarity.
The author received no funding for this work and declares no competing
interests.

\section*{AI assistance}

OpenAI's GPT-5.6 Pro through ChatGPT and Aristotle~\citep{achim-et-al-2025}
were used in proof development, adversarial checking, editorial
restructuring, and the Lean formalisation.  In particular, ChatGPT reviews
informed the checking and presentation of the local forest-counting and
capacity arguments in Section~10; Aristotle supplied the abstract
incidence-forest construction distinguished from fixed-atom realization in
the verification section.  The author takes full responsibility for the
arguments, citations, and final manuscript; newly incorporated material in
this review draft remains subject to the author's review.

\begin{thebibliography}{99}
\interlinepenalty=10000
\providecommand{\natexlab}[1]{#1}
\providecommand{\url}[1]{\texttt{#1}}
\expandafter\ifx\csname urlstyle\endcsname\relax
  \providecommand{\doi}[1]{doi: #1}\else
  \providecommand{\doi}{doi: \begingroup \urlstyle{rm}\Url}\fi

\bibitem[Erd{\H{o}}s(1995)]{erdos1995}
Paul Erd{\H{o}}s.
\newblock On some problems in combinatorial set theory.
\newblock \emph{Publikacije Instituta Matemati\v{c}kog (Beograd) (N.S.)},
  57(71):\penalty0 61--65, 1995.

\bibitem[Bloom(2026)]{bloom593}
Thomas~F. Bloom.
\newblock Erd{\H{o}}s problem 593.
\newblock \url{https://www.erdosproblems.com/593}, 2026.
\newblock Accessed 11 July 2026.

\bibitem[Li(2026)]{li2026}
Eric Li.
\newblock A resolution of {Erd{\H{o}}s Problems 593 and 1177}: Obligatory
  triple systems and exact spectra.
\newblock arXiv:2606.24882, submitted 23 June 2026.
\newblock \doi{10.48550/arXiv.2606.24882}.
\newblock URL \url{https://arxiv.org/abs/2606.24882}.

\bibitem[Whitney(1932)]{whitney1932}
Hassler Whitney.
\newblock Non-separable and planar graphs.
\newblock \emph{Trans. Amer. Math. Soc.}, 34\penalty0 (2):\penalty0
  339--362, 1932.
\newblock \doi{10.1090/S0002-9947-1932-1501641-7}.

\bibitem[Diestel(2017)]{diestel2017}
Reinhard Diestel.
\newblock \emph{Graph Theory}, fifth edition.
\newblock Graduate Texts in Mathematics, volume~173. Springer, Berlin,
  2017.
\newblock \doi{10.1007/978-3-662-53622-3}.

\bibitem[Stanley(2012)]{stanley2012}
Richard~P. Stanley.
\newblock \emph{Enumerative Combinatorics}, volume~1, second edition.
\newblock Cambridge Studies in Advanced Mathematics, volume~49. Cambridge
  University Press, Cambridge, 2012.
\newblock \doi{10.1017/CBO9781139058520}.

\bibitem[Comtet(1974)]{comtet1974}
Louis Comtet.
\newblock \emph{Advanced Combinatorics: The Art of Finite and Infinite Expansions}.
\newblock D. Reidel, Dordrecht, 1974.
\newblock \doi{10.1007/978-94-010-2196-8}.

\bibitem[Csikv{\'a}ri and Lin(2017)]{csikvarilin2017}
P{\'e}ter Csikv{\'a}ri and Zhicong Lin.
\newblock Sidorenko's conjecture, colorings and independent sets.
\newblock \emph{Electron. J. Combin.}, 24(1):P1.2, 2017.
\newblock \doi{10.37236/6019}.
\newblock URL \url{https://arxiv.org/abs/1603.05888}.

\bibitem[Szegedy(2015)]{szegedy2015}
Bal{\'a}zs Szegedy.
\newblock Sparse graph limits, entropy maximization and transitive graphs.
\newblock arXiv:1504.00858, 2015.
\newblock \doi{10.48550/arXiv.1504.00858}.
\newblock URL \url{https://arxiv.org/abs/1504.00858}.

\bibitem[Lov{\'a}sz(2010)]{lovasz2010}
L{\'a}szl{\'o} Lov{\'a}sz.
\newblock Subgraph densities in signed graphons and the local Sidorenko conjecture.
\newblock arXiv:1004.3026, 2010.
\newblock URL \url{https://arxiv.org/abs/1004.3026}.

\bibitem[Sidorenko(2021)]{sidorenko2021}
Alexander Sidorenko.
\newblock Inequalities for doubly nonnegative functions.
\newblock \emph{Electron. J. Combin.}, 28(1):P1.32, 2021.
\newblock \doi{10.37236/8947}.
\newblock URL \url{https://arxiv.org/abs/1905.08210}.

\bibitem[Conlon et~al.(2012)Conlon, Fox, and Sudakov]{conlonfoxsudakov2012}
David Conlon, Jacob Fox, and Benny Sudakov.
\newblock Sidorenko's conjecture for a class of graphs: an exposition.
\newblock arXiv:1209.0184, 2012.
\newblock URL \url{https://arxiv.org/abs/1209.0184}.

\bibitem[Im et~al.(2026)Im, Li, and Liu]{imliliu2026}
Seonghyuk Im, Ruonan Li, and Hong Liu.
\newblock Sidorenko's conjecture for subdivisions and theta substitutions.
\newblock \emph{Combin. Probab. Comput.}, 35(2):269--279, 2026.
\newblock \doi{10.1017/S0963548325100242}.

\bibitem[Kostochka et~al.(2015)Kostochka, Mubayi, and Verstra{\"e}te]{kostochkamubayiverstraete2015}
Alexandr Kostochka, Dhruv Mubayi, and Jacques Verstra{\"e}te.
\newblock Tur{\'a}n problems and shadows III: Expansions of graphs.
\newblock \emph{SIAM J. Discrete Math.}, 29\penalty0 (2):\penalty0 868--876, 2015.
\newblock \doi{10.1137/140977138}.

\bibitem[Erd{\H{o}}s and Hajnal(1966)]{erdoshajnal1966}
Paul Erd{\H{o}}s and Andr{\'a}s Hajnal.
\newblock On chromatic number of graphs and set-systems.
\newblock \emph{Acta Math. Acad. Sci. Hungar.}, 17\penalty0
  (1--2):\penalty0 61--99, 1966.
\newblock \doi{10.1007/BF02020444}.

\bibitem[Erd{\H{o}}s et~al.(1973)Erd{\H{o}}s, Hajnal, and Rothschild]{ehr1973}
Paul Erd{\H{o}}s, Andr{\'a}s Hajnal, and Bruce~L. Rothschild.
\newblock On chromatic number of graphs and set-systems.
\newblock In \emph{Cambridge Summer School in Mathematical Logic (Cambridge,
  1971)}, volume 337 of \emph{Lecture Notes in Mathematics}, pages 531--538.
  Springer, Berlin and New York, 1973.

\bibitem[Erd{\H{o}}s et~al.(1975)Erd{\H{o}}s, Galvin, and Hajnal]{egh1975}
Paul Erd{\H{o}}s, Fred Galvin, and Andr{\'a}s Hajnal.
\newblock On set-systems having large chromatic number and not containing
  prescribed subsystems.
\newblock In \emph{Infinite and Finite Sets (Keszthely, 1973)}, volume~10 of
  \emph{Colloquia Mathematica Societatis J{\'a}nos Bolyai}, pages 425--513.
  North-Holland, 1975.

\bibitem[Komj{\'a}th(2001)]{komjath2001}
P{\'e}ter Komj{\'a}th.
\newblock Some remarks on obligatory subsystems of uncountably chromatic triple
  systems.
\newblock \emph{Combinatorica}, 21\penalty0 (2):\penalty0 233--238, 2001.
\newblock \doi{10.1007/s004930100021}.

\bibitem[Hajnal and Komj{\'a}th(2008)]{hajnalkomjath2008}
Andr{\'a}s Hajnal and P{\'e}ter Komj{\'a}th.
\newblock Obligatory subsystems of triple systems.
\newblock \emph{Acta Math. Hungar.}, 119\penalty0 (1--2):\penalty0
  1--13, 2008.
\newblock \doi{10.1007/s10474-007-6231-2}.

\bibitem[Komj{\'a}th(2008)]{komjath2008}
P{\'e}ter Komj{\'a}th.
\newblock An uncountably chromatic triple system.
\newblock \emph{Acta Math. Hungar.}, 121\penalty0 (1--2):\penalty0
  79--92, 2008.
\newblock \doi{10.1007/s10474-008-7179-6}.

\bibitem[Reiher(2024)]{reiher2024}
Christian Reiher.
\newblock Obligatory hypergraphs.
\newblock arXiv:2403.11223, 2024.
\newblock \doi{10.1090/proc/17021}.
\newblock URL \url{https://arxiv.org/abs/2403.11223}.

\bibitem[Erd{\H{o}}s(1959)]{erdos1959}
Paul Erd{\H{o}}s.
\newblock Graph theory and probability.
\newblock \emph{Canad. J. Math.}, 11:\penalty0 34--38, 1959.
\newblock \doi{10.4153/CJM-1959-003-9}.

\bibitem[Erd{\H{o}}s and Lov{\'a}sz(1975)]{erdoslovasz1975}
Paul Erd{\H{o}}s and L{\'a}szl{\'o} Lov{\'a}sz.
\newblock Problems and results on 3-chromatic hypergraphs and some related questions.
\newblock In Andr{\'a}s Hajnal, Richard Rado, and Vera~T. S{\'o}s, editors,
  \emph{Infinite and Finite Sets}, volume~10 of \emph{Colloquia Mathematica
  Societatis J{\'a}nos Bolyai}, pages 609--627. North-Holland, Amsterdam, 1975.
\newblock URL \url{https://www.renyi.hu/~p_erdos/1975-34.pdf}.

\bibitem[Kostochka and R{\"o}dl(2010)]{kostochkarodl2010}
A.~V. Kostochka and V.~R{\"o}dl.
\newblock Constructions of sparse uniform hypergraphs with high chromatic number.
\newblock \emph{Random Structures Algorithms}, 36\penalty0 (1):\penalty0 46--56, 2010.
\newblock \doi{10.1002/rsa.20293}.

\bibitem[Alon et~al.(2016)Alon, Kostochka, Reiniger, West, and Zhu]{alonetal2016}
Noga Alon, Alexandr Kostochka, Benjamin Reiniger, Douglas~B. West, and Xuding Zhu.
\newblock Coloring, sparseness and girth.
\newblock \emph{Israel J. Math.}, 214\penalty0 (1):\penalty0 315--331, 2016.
\newblock \doi{10.1007/s11856-016-1361-2}.

\bibitem[Wang et~al.(2026)Wang, Duan, Gerbner, and Hama Karim]{wang2026}
Yichen Wang, Mengyu Duan, D{\'a}niel Gerbner, and Hilal Hama Karim.
\newblock On the largest chromatic number of {$F$}-free hypergraphs.
\newblock arXiv:2604.21551, 2026.
\newblock \doi{10.48550/arXiv.2604.21551}.
\newblock URL \url{https://arxiv.org/abs/2604.21551}.

\bibitem[de~Bruijn and Erd{\H{o}}s(1951)]{debruijn1951}
Nicolaas~Govert de~Bruijn and Paul Erd{\H{o}}s.
\newblock A colour problem for infinite graphs and a problem in the theory of
  relations.
\newblock \emph{Nederl. Akad. Wetensch. Proc. Ser. A}, 54\penalty0
  (5):\penalty0 371--373, 1951.
\newblock \doi{10.1016/S1385-7258(51)50053-7}.

\bibitem[Erd{\H{o}}s and Rado(1956)]{erdosrado1956}
Paul Erd{\H{o}}s and Richard Rado.
\newblock A partition calculus in set theory.
\newblock \emph{Bull. Amer. Math. Soc.}, 62\penalty0
  (5):\penalty0 427--489, 1956.
\newblock \doi{10.1090/S0002-9904-1956-10036-0}.

\bibitem[Erd{\H{o}}s et~al.(1984)Erd{\H{o}}s, Hajnal, M{\'a}t{\'e}, and Rado]{ehmr1984}
Paul Erd{\H{o}}s, Andr{\'a}s Hajnal, Attila M{\'a}t{\'e}, and Richard Rado.
\newblock \emph{Combinatorial Set Theory: Partition Relations for Cardinals}.
\newblock Studies in Logic and the Foundations of Mathematics, volume~106.
  North-Holland, Amsterdam, and Akad{\'e}miai Kiad{\'o}, Budapest, 1984.

\bibitem[Bahmanian and \v{S}ajna(2015)]{bahmanian2015}
M.~Amin Bahmanian and Mateja \v{S}ajna.
\newblock Connection and separation in hypergraphs.
\newblock \emph{Theory Appl. Graphs}, 2\penalty0 (2):\penalty0 Article~5, 2015.
\newblock \doi{10.20429/tag.2015.020205}.
\newblock URL \url{https://arxiv.org/abs/1504.04274}.

\bibitem[Tarjan(1974)]{tarjan1974}
R.~Endre Tarjan.
\newblock A note on finding the bridges of a graph.
\newblock \emph{Inform. Process. Lett.}, 2\penalty0 (6):\penalty0 160--161, 1974.
\newblock \doi{10.1016/0020-0190(74)90003-9}.

\bibitem[Tarjan(1972)]{tarjan1972}
Robert Tarjan.
\newblock Depth-first search and linear graph algorithms.
\newblock \emph{SIAM J. Comput.}, 1\penalty0 (2):\penalty0 146--160, 1972.
\newblock \doi{10.1137/0201010}.

\bibitem[Simon et~al.(2011)Simon, Tittmann, and Trinks]{simon2011}
Frank Simon, Peter Tittmann, and Martin Trinks.
\newblock Counting connected set partitions of graphs.
\newblock \emph{Electron. J. Combin.}, 18\penalty0 (1):\penalty0 P14, 2011.
\newblock \doi{10.37236/501}.
\newblock URL \url{https://arxiv.org/abs/1005.1726}.

\bibitem[Stanley(2007)]{stanley2007}
Richard~P. Stanley.
\newblock An introduction to hyperplane arrangements.
\newblock In Ezra Miller, Victor Reiner, and Bernd Sturmfels, editors,
  \emph{Geometric Combinatorics}, volume~13 of \emph{IAS/Park City
  Mathematics Series}, pages 389--496. American Mathematical Society,
  Providence, RI, 2007.
\newblock URL \url{https://math.mit.edu/~rstan/arrangements/arr.html}.

\bibitem[Gr{\"a}tzer(2011)]{gratzer2011}
George Gr{\"a}tzer.
\newblock \emph{Lattice Theory: Foundation}.
\newblock Birkh{\"a}user, Basel, 2011.
\newblock \doi{10.1007/978-3-0348-0018-1}.

\bibitem[Rota(1964)]{rota1964}
Gian-Carlo Rota.
\newblock On the foundations of combinatorial theory. I. Theory of
  M\"obius functions.
\newblock \emph{Z. Wahrscheinlichkeitstheorie verw. Geb.}, 2:\penalty0
  340--368, 1964.
\newblock \doi{10.1007/BF00531932}.

\bibitem[Schmitt(1994)]{schmitt1994}
William~R. Schmitt.
\newblock Incidence Hopf algebras.
\newblock \emph{J. Pure Appl. Algebra}, 96\penalty0 (3):\penalty0 299--330, 1994.
\newblock \doi{10.1016/0022-4049(94)90105-8}.

\bibitem[Lieb(1968)]{lieb1968}
Elliott~H. Lieb.
\newblock Concavity properties and a generating function for Stirling numbers.
\newblock \emph{J. Combin. Theory}, 5\penalty0 (2):\penalty0 203--206, 1968.
\newblock \doi{10.1016/S0021-9800(68)80057-2}.

\bibitem[Pitman(1997)]{pitman1997}
Jim Pitman.
\newblock Probabilistic bounds on the coefficients of polynomials with only real zeros.
\newblock \emph{J. Combin. Theory Ser. A}, 77\penalty0 (2):\penalty0 279--303, 1997.
\newblock \doi{10.1006/jcta.1997.2747}.

\bibitem[de~Moura and Ullrich(2021)]{demoura2021}
Leonardo de~Moura and Sebastian Ullrich.
\newblock The Lean~4 theorem prover and programming language.
\newblock In Andr\'e Platzer and Geoff Sutcliffe, editors,
  \emph{Automated Deduction---CADE~28}, volume~12699 of \emph{Lecture
  Notes in Computer Science}, pages 625--635. Springer, Cham, 2021.
\newblock \doi{10.1007/978-3-030-79876-5_37}.

\bibitem[The mathlib Community(2020)]{mathlib2020}
The mathlib Community.
\newblock The Lean mathematical library.
\newblock In \emph{Proceedings of the 9th ACM SIGPLAN International
  Conference on Certified Programs and Proofs}, pages 367--381.
  Association for Computing Machinery, New York, 2020.
\newblock \doi{10.1145/3372885.3373824}.

\bibitem[Avigad and Harrison(2014)]{avigadharrison2014}
Jeremy Avigad and John Harrison.
\newblock Formally verified mathematics.
\newblock \emph{Commun. ACM}, 57\penalty0 (4):\penalty0 66--75, 2014.
\newblock \doi{10.1145/2591012}.

\bibitem[Carneiro(2025)]{carneiro2025}
Mario Carneiro.
\newblock Lean4Lean: Verifying a typechecker for Lean, in Lean.
\newblock arXiv:2403.14064v3, revised 14 September 2025.
\newblock URL \url{https://arxiv.org/abs/2403.14064v3}.

\bibitem[Achim et~al.(2025)Achim, Best, Der, F{\'e}d{\'e}rico, Gukov,
  Halpern-Leistner, Henningsgard, Kudryashov, Meiburg, Michelsen, Patterson,
  Rodriguez, Scharff, Shanker, Sicca, Sowrirajan, Swope, Tamas, Tenev, Thomm,
  Williams, and Wu]{achim-et-al-2025}
Tudor Achim, Alex Best, Kevin Der, Math{\"i}s F{\'e}d{\'e}rico, Sergei Gukov,
  Daniel Halpern-Leistner, Kirsten Henningsgard, Yury Kudryashov, Alexander
  Meiburg, Martin Michelsen, Riley Patterson, Eric Rodriguez, Laura Scharff,
  Vikram Shanker, Vladimir Sicca, Hari Sowrirajan, Aidan Swope, Matyas Tamas,
  Vlad Tenev, Jonathan Thomm, Harold Williams, and Lawrence Wu.
\newblock {Aristotle}: {IMO}-level automated theorem proving, 2025.
\newblock URL \url{https://arxiv.org/abs/2510.01346}.
\newblock arXiv:2510.01346.

\end{thebibliography}

\end{document}
